\documentclass[12pt,a4paper]{report}

% ===== Packages =====
\usepackage[top=2.5cm, bottom=2.5cm, left=2.5cm, right=2.5cm,headheight=2.0cm,headsep=0.5cm,footskip=1.2cm]{geometry}
\usepackage[table,dvipsnames,svgnames,x11names]{xcolor}
\usepackage{graphicx}
\usepackage{amsmath,amssymb}
\usepackage{fancyhdr}
\usepackage{titlesec}
\usepackage{enumitem}
\usepackage{float}
\usepackage{placeins}
\usepackage{needspace}
\usepackage[T1]{fontenc}
\usepackage[utf8]{inputenc}
\usepackage[french]{babel}
% Disable French automatic spacing before high punctuation (: ! ?)
\frenchsetup{AutoSpacePunctuation=false}
\usepackage{newtxtext,newtxmath}
\usepackage{microtype}
\usepackage{setspace}
\usepackage{tikz}
\usepackage{array}
\usepackage{tabularx}
\usepackage{caption}
\usepackage{tocloft}
\usepackage{longtable}
\usepackage{url}




% ===== Chapter-based numbering =====
\setcounter{secnumdepth}{3}
\setcounter{tocdepth}{3}

% EMSI PFE typography (Times New Roman via newtx):
% - Body: 12pt, interligne 1,5
% - Front-matter titles: 24pt bold centered
% - Chapter titles: 18pt bold
% - Section (x.1): 16pt bold
% - Subsection (x.1.1): 14pt bold
% - Captions: bold centered

% Chapter format
\titleformat{\chapter}
  {\normalfont\fontsize{18}{22}\selectfont\bfseries\color{titlegreen}}
  {Chapitre \thechapter.}{1em}{}

% Section numbering
\renewcommand{\thesection}{\thechapter.\arabic{section}}
\renewcommand{\thesubsection}{\thesection.\arabic{subsection}}
\renewcommand{\thesubsubsection}{\thesubsection.\arabic{subsubsection}}

% Figures & Tables numbering
\numberwithin{figure}{chapter}
\numberwithin{table}{chapter}

\renewcommand{\thefigure}{\thechapter.\arabic{figure}}
\renewcommand{\thetable}{\thechapter.\arabic{table}}


% ===== Colors =====
\definecolor{titlegreen}{RGB}{0,176,80}        % Main heading green
\definecolor{tableheadergreen}{RGB}{0,176,80}  % Glossary highlighted rows
\definecolor{darkgreen}{RGB}{0,128,0}          % Slightly darker green for chapter rules

% ===== Section formatting (EMSI sizes, Times New Roman) =====
\titleformat{\section}
  {\fontsize{16}{20}\selectfont\bfseries\color{titlegreen}}
  {\thesection.}{0.6em}{}
\titleformat{\subsection}
  {\fontsize{14}{18}\selectfont\bfseries\color{titlegreen}}
  {\thesubsection}{0.8em}{}
\titleformat{\subsubsection}
  {\fontsize{12}{15}\selectfont\bfseries\color{titlegreen}}
  {\thesubsubsection}{0.5em}{}

\titlespacing*{\section}{0pt}{14pt}{8pt}
\titlespacing*{\subsection}{1.5em}{12pt}{6pt}
\titlespacing*{\subsubsection}{3em}{10pt}{4pt}

\setlength{\parindent}{0pt}     % no paragraph indent
\setlength{\parskip}{0.5em}     % space between paragraphs

% Caption setup (EMSI: gras, centré)
\captionsetup{font=bf,labelfont=bf,justification=centering,labelsep=period}
\captionsetup[figure]{name=Figure}
\captionsetup[table]{name=Tableau,position=top}
\renewcommand{\figurename}{Figure}
\renewcommand{\tablename}{Tableau}

% ===== Header / Footer =====
\pagestyle{fancy}
\fancyhf{}
\renewcommand{\headrulewidth}{0pt}
\renewcommand{\footrulewidth}{0pt}

% Safe image command (falls back to a styled label if image missing)
\newcommand{\safeimage}[2]{%
  \IfFileExists{#2}{\includegraphics[#1]{#2}}{%
    \fbox{\scriptsize Missing: \detokenize{#2}}%
  }%
}

% Header logos
\fancyhead[R]{\safeimage{height=1.2cm}{images/tripass_logo.jpg}}
\fancyhead[L]{\safeimage{height=1.2cm}{images/emsi_logo.jpeg}}

% Footer page number with horizontal rules and curly braces
\fancyfoot[C]{%
  \rule[0.4ex]{0.32\textwidth}{0.5pt}%
  \hspace{4pt}\Large$\bigg($\normalsize\,\thepage\,\Large$\bigg)$\normalsize\hspace{4pt}%
  \rule[0.4ex]{0.32\textwidth}{0.5pt}%
}

% Force plain pages (LoF, LoT, ToC) to use fancy header (logos stay visible)
\fancypagestyle{plain}{
  \fancyhf{}
  \fancyhead[R]{\safeimage{height=1.5cm}{images/tripass_logo.jpg}}
  \fancyhead[L]{\safeimage{height=1.5cm}{images/emsi_logo.jpeg}}
  \fancyfoot[C]{%
    \rule[0.4ex]{0.32\textwidth}{0.5pt}%
    \hspace{4pt}\Large$\bigg($\normalsize\,\thepage\,\Large$\bigg)$\normalsize\hspace{4pt}%
    \rule[0.4ex]{0.32\textwidth}{0.5pt}%
  }
  \renewcommand{\headrulewidth}{0pt}
}

% Custom heading for front-matter titles (EMSI: Times New Roman 24pt bold centered)
\newcommand{\centeredgreentitle}[1]{%
  \vspace*{0.8cm}%
  \begin{center}
    {\fontsize{24}{28}\selectfont\bfseries\color{titlegreen} #1}
  \end{center}
  \vspace{0.5cm}
}

% Chapter title page: "Chapitre X :" top-left, title centered between rules
\newcommand{\chapterpagetitle}[1]{%
  \vspace*{1.2cm}%
  {\noindent\fontsize{18}{22}\selectfont\bfseries\color{titlegreen} Chapitre \thechapter :}\par
  \vspace*{6.0cm}%
  \noindent\textcolor{darkgreen}{\rule{\textwidth}{2pt}}%
  \vspace{1.5cm}%
  \begin{center}
    {\fontsize{24}{28}\selectfont\bfseries\color{titlegreen} #1}%
  \end{center}
  \vspace{1.5cm}%
  \noindent\textcolor{darkgreen}{\rule{\textwidth}{2pt}}%
  \vfill
  \newpage
}

% TOC entry helpers
\newcommand{\addfrontmattertotoc}[1]{\addcontentsline{toc}{chapter}{#1}}
\newcommand{\addintrototoc}{\addcontentsline{toc}{chapter}{Introduction générale}}




% spacing between title and entries
\setlength{\cftbeforetoctitleskip}{0pt}
\setlength{\cftaftertoctitleskip}{10pt}
\setlength{\cftbeforeloftitleskip}{0pt}
\setlength{\cftafterloftitleskip}{10pt}
\setlength{\cftbeforelottitleskip}{0pt}
\setlength{\cftafterlottitleskip}{10pt}

% dots style
\renewcommand{\cftdotsep}{1}
\renewcommand{\cftchapdotsep}{\cftdotsep}

% figure number style
\renewcommand{\cftfigpresnum}{Figure~}
\renewcommand{\cftfigaftersnum}{: }
\setlength{\cftfignumwidth}{2cm}

% table number style
\renewcommand{\cfttabpresnum}{Tableau~}
\renewcommand{\cfttabaftersnum}{: }
\setlength{\cfttabnumwidth}{2.2cm}

% indentation levels


\setlength{\cftsecindent}{1.5em}
\setlength{\cftsubsecindent}{3em}
\setlength{\cftsubsubsecindent}{4.5em}

% ===== Document =====
\begin{document}
\onehalfspacing
\pagenumbering{roman}
% ===============================================================
% PAGE 1 — DÉDICACE
% ===============================================================
\centeredgreentitle{Dédicace}
\addfrontmattertotoc{Dédicace}

\vspace{0.5cm}

\begin{center}
\fontsize{16}{22}\selectfont\itshape
Tout d’abord je voudrais remercier le dieu tout puissant et miséricordieux qui m’a donné la force, l’intelligence et la patience pour accomplir ce travail.

A mes chers parents

Aucune dédicace ne saurait exprimer mon respect, mon amour éternel et ma considération pour les sacrifices que vous avez consentie pour mon instruction et mon bien-être.

Je vous remercie pour tout le soutien et l’amour que vous m’avez donné et j’espère que votre bénédiction m’accompagne toujours.

Que ce modeste travail soit l’exaucement de vos vœux tant formulés, le fruit de vos innombrables sacrifices, bien qu’on ne vous en acquitte jamais assez.

Puisse Dieu, le Très Haut, vous accorde santé, bonheur et longue vie et faire en sorte que jamais nous ne vous déçoive.
\end{center}

\newpage

% ===============================================================
% PAGE 2 — REMERCIEMENTS
% ===============================================================
\centeredgreentitle{Remerciements}
\addfrontmattertotoc{Remerciements}

\vspace{0.5cm}

\begin{center}
\fontsize{16}{22}\selectfont\itshape
Au terme de ce travail, je tiens à exprimer ma profonde gratitude à toutes les personnes qui ont contribué, de près ou de loin, à la réalisation de ce projet.

Je remercie en premier lieu mon encadrant pédagogique, Mme. Mariem Malaine, pour son accompagnement, ses conseils pertinents ainsi que le suivi régulier tout au long de cette période.

Mes remerciements s’adressent également à Mr. Abid Khirani, Directeur Général de Tripass, pour m’avoir accueilli au sein de l’entreprise et pour la confiance qu’il m’a accordée dès mon arrivée.

Je remercie également Mr. Amine Ahbouch, responsable de la direction développement \& technique, pour son encadrement technique tout au long du stage, pour son soutien et sa coopération durant mon intégration au sein de l’entreprise.

Je n’oublie pas l’ensemble de l’équipe Tripass pour leur esprit d’équipe, leur disponibilité et leur contribution à la réussite de cette expérience professionnelle enrichissante.

Enfin, je tiens à remercier les membres du jury pour l’intérêt qu’ils portent à mon travail en acceptant de l’évaluer.
\end{center}

\newpage

% ===============================================================
% PAGE 3 — ABSTRACT
% ===============================================================
\centeredgreentitle{Abstract}
\addfrontmattertotoc{Abstract}

\vspace{0.5cm}

{\fontsize{14}{20}\selectfont
To complete my computer science education and begin my professional career, I undertook a 4-month internship dedicated to developing \textit{Agent Market}, a web platform for managing and optimizing digital marketing campaigns.

My work consisted in designing and building a full-stack application for campaign management and automation, using Next.js and React for interactive, responsive interfaces and an optimized user experience.

Following a Waterfall approach, I carried out a needs analysis and a UML design, then integrated MongoDB, Python for intelligent processing, Redis and Celery for background tasks, and Docker Compose for containerization and deployment. Testing throughout development ensured reliability, performance, and scalability, and I contributed to deploying the application in a production-like environment.

This internship strengthened my skills in full-stack development, distributed systems, and project management, preparing me for the challenges of modern software engineering.
\par}

\vspace{0.5cm}

\begin{center}
{\sffamily\fontsize{14}{18}\selectfont\itshape Keywords : Web Application, Full Stack Development, Database, UML, Next.js, React, Fastify, MongoDB, Python, Redis, Celery, Docker Compose, Waterfall}
\end{center}
\newpage
% ===============================================================
% PAGE 4 — RÉSUMÉ
% ===============================================================
\centeredgreentitle{Résumé}
\addfrontmattertotoc{Résumé}

\vspace{0.5cm}

{\fontsize{14}{20}\selectfont
Afin de compléter ma formation en informatique et de débuter ma carrière professionnelle, j'ai effectué un stage de quatre mois dédié au développement d'\textit{Agent Market}, une plateforme web de gestion et d'optimisation des campagnes de marketing digital.

Mon travail a consisté à concevoir et réaliser une application full-stack pour la gestion et l'automatisation des campagnes, en utilisant Next.js et React pour des interfaces interactives, responsives et une expérience utilisateur optimisée.

Suivant une approche Waterfall, j'ai mené une analyse des besoins et une conception UML, puis intégré MongoDB, Python pour le traitement intelligent, Redis et Celery pour les tâches en arrière-plan, ainsi que Docker Compose pour la conteneurisation et le déploiement. Des tests tout au long du développement ont assuré la fiabilité, la performance et la scalabilité, et j'ai contribué au déploiement de l'application dans un environnement proche de la production.

Ce stage m'a permis de renforcer mes compétences en développement full-stack, en systèmes distribués et en gestion de projet, me préparant aux défis de l'ingénierie logicielle moderne.
\par}

\vspace{0.8cm}

\begin{center}
{\sffamily\fontsize{14}{18}\selectfont\itshape Mots clés : Application Web, Développement Full Stack, Base de données,\\
UML, Next.js, React, Fastify, MongoDB, Python, Redis, Celery, Docker\\
Compose, Waterfall}
\end{center}

\newpage

% ===============================================================
% PAGE 5 — GLOSSAIRE
% ===============================================================
\centeredgreentitle{Glossaire}
\addfrontmattertotoc{Glossaire}

\vspace{0.5cm}

\textbf{API} : Application Programming Interface.

\textbf{EMSI} : Ecole Marocaine des Sciences de l'Ingénieur.

\textbf{GIT} : Global Information Tracker.

\textbf{MIAGE} : Méthodes Informatique Appliquées à la Gestion d'Entreprise.

\textbf{MVC} : Model-View-Controller.

\textbf{UML} : Unified Modeling Language.

\newpage

% ===============================================================
% PAGE 6 — TABLE OF FIGURES
% ===============================================================
\centeredgreentitle{Liste des figures}
\addfrontmattertotoc{Liste des figures}


\renewcommand{\listfigurename}{}
\listoffigures

\newpage

% ===============================================================
% PAGE 7 — INDEX OF TABLES
% ===============================================================
\centeredgreentitle{Liste des tableaux}
\addfrontmattertotoc{Liste des tableaux}

\renewcommand{\listtablename}{}
\listoftables

\newpage

% ===============================================================
% PAGE 8 — TABLE OF CONTENTS (part 1)
% ===============================================================
\centeredgreentitle{Table des matières}
\addfrontmattertotoc{Table des matières}

\renewcommand{\contentsname}{}
\tableofcontents

\newpage
\pagenumbering{arabic}

% ===============================================================
% PAGE 9 — INTRODUCTION GÉNÉRALE
% ===============================================================
\centeredgreentitle{Introduction générale}
\addintrototoc

L'intelligence artificielle (IA), en tant que technologie avancée de traitement et d'analyse automatisée de l'information, occupe aujourd'hui une place centrale dans de nombreux domaines, qu'ils soient scientifiques, professionnels, privés ou publics. Grâce à l'évolution rapide des technologies numériques telles que le machine learning, le traitement du langage naturel et l'analyse prédictive, les organisations disposent désormais d'outils performants pour exploiter efficacement leurs données. Ces technologies permettent d'automatiser des tâches complexes, d'améliorer la qualité des analyses et de soutenir les processus décisionnels dans un environnement numérique en constante transformation.

\vspace{0.3cm}

Dans le domaine du marketing digital, l'intégration de ces technologies contribue fortement à l'amélioration des performances des entreprises. La multiplication des plateformes numériques, des outils publicitaires et des solutions d'analyse permet aux organisations d'optimiser leurs campagnes, de cibler plus précisément leurs audiences et de mesurer les résultats en temps réel. Cette capacité d'analyse et d'adaptation rapide favorise une meilleure productivité, une augmentation du retour sur investissement et une amélioration globale de l'efficacité des stratégies marketing.

\vspace{0.3cm}

Cependant, malgré ces avancées, certaines limites persistent dans la gestion marketing multi-plateformes. L'utilisation simultanée de plusieurs outils indépendants peut entraîner une dispersion des données, une complexité accrue dans la gestion des campagnes et une difficulté à obtenir une vision globale des performances. De plus, le manque de centralisation et d'automatisation peut ralentir la prise de décision et réduire l'efficacité des équipes marketing, notamment dans un environnement où la réactivité est essentielle.

\vspace{0.3cm}


Dans ce contexte, ce rapport est structuré en plusieurs chapitres. Le premier chapitre présente l’organisme d’accueil ainsi que le contexte général du projet. Le deuxième chapitre est consacré à l’analyse du projet, incluant l’analyse des besoins et la présentation des fonctionnalités de l’assistant IA. Le troisième chapitre aborde l’étude conceptuelle à travers la modélisation UML. Le quatrième chapitre présente les outils et technologies utilisés pour la réalisation du projet. Enfin, le cinquième chapitre est dédié à la réalisation de la solution, en présentant les interfaces de l’application accompagnées de leurs explications.

\newpage

% ===============================================================
% PAGE 10 — CHAPTER 1 TITLE PAGE
% ===============================================================



\refstepcounter{chapter}

\addcontentsline{toc}{chapter}{Chapitre \thechapter : Organisme d'accueil et Contexte général du Projet}

\chapterpagetitle{Organisme d'accueil et Contexte général du Projet}


% ===============================================================
% PAGE 11 — CHAPTER CONTENT (Introduction + Tripass SARL)
% ===============================================================

\section{Introduction}

La présentation du contexte général du projet a pour but de situer le projet dans son environnement organisationnel et contextuel. Ce chapitre introduit donc l'organisme d'accueil, décrit la problématique à traiter et présente la démarche suivie pour la réalisation du projet.

\section{Presentation generale d'organisme d'accueil :}

\subsection{Tripass SARL}

\textbf{Tripass SARL} est une startup marocaine fondée le \textbf{20/12/2019} et spécialisée dans la digitalisation du transport. Elle s'est particulièrement distinguée par sa solution innovante \textbf{WeTaxi}, un service dédié à la gestion simplifiée des taxis, notamment dans plusieurs grands aéroports du Maroc. Tripass contribue à moderniser l'expérience de transport en proposant une approche plus fluide, plus rapide et plus transparente pour les voyageurs. Son objectif principal est de faciliter l'accès aux taxis tout en améliorant l'organisation des guichets et la qualité du service offert aux clients.

Comme illustré dans la Figure~\ref{fig:tripass_logo}, le logo de Tripass représente l'identité visuelle de l'entreprise et reflète son positionnement dans le secteur du transport digitalisé.

\begin{figure}[H]
\centering
\safeimage{width=0.35\textwidth}{images/tripass_logo.jpg}
\caption{Logo Tripass}
\label{fig:tripass_logo}
\end{figure}

L'agence se démarque par ses solutions numériques orientées vers le transport et l'accueil des passagers, avec une attention particulière portée à l'efficacité, à la simplicité d'utilisation et à la satisfaction des usagers.

\subsection{Fiche Technique}

Dans le cadre de la présentation générale de l'entreprise, il est nécessaire de fournir une fiche technique regroupant les principales informations administratives et organisationnelles de la société Tripass. Cette fiche permet de présenter de manière synthétique l'identité de l'entreprise, son domaine d'activité, ses coordonnées ainsi que ses informations légales. Elle constitue ainsi une référence essentielle pour mieux comprendre le cadre institutionnel et organisationnel de l'entreprise.


\vspace{0.3cm}

Le tableau suivant présente la fiche technique de la société Tripass :

\vspace{0.5cm}

% ===================== TABLE =====================
\begin{table}[H]
\centering
\caption{Fiche Technique de Tripass}

\renewcommand{\arraystretch}{1.6}
\begin{tabular}{|p{0.28\textwidth}|p{0.62\textwidth}|}
\hline
\textbf{Élément} & \textbf{Détail} \\
\hline
Dénomination & Tripass SARL \\
\hline
Date de création & 20/12/2019 \\
\hline
Forme juridique & Societe A Responsabilite Limitee. \\
\hline
Siège social & Tourisme \\
\hline
Effectif & 93 employés \\
\hline
Activité & Tourisme \\
\hline
Clients cibles & Touristes \\
\hline
RC & 101243 \\
\hline
ICE & 002371582000090 \\
\hline
CNSS & 1840113 \\
\hline
Téléphone société & 06 88 99 94 06 \\
\hline
Téléphone fixe & 05 24 43 09 41 \\
\hline
Email & ak@casky.io \\
\hline
\end{tabular}

\renewcommand{\arraystretch}{1.0}
\end{table}

\vspace{0.5cm}

% ===================== NEXT SECTION =====================
\subsection{Rôles de l'entreprise}

\hspace{1.2em}Tripass joue un rôle important dans la modernisation du transport au Maroc, en particulier à travers sa solution \textbf{WeTaxi}. Pour de nombreux voyageurs, le premier contact avec le pays se fait par un trajet en taxi aéroportuaire. Or, ce segment a longtemps été marqué par l’absence de tarification standardisée et des négociations à la descente d’avion. WeTaxi digitalise le taxi aéroportuaire en reliant les aéroports du Royaume aux centres-villes selon un modèle prévisible, transparent et sécurisé. Le service est déployé dans plusieurs villes, notamment \textbf{Casablanca}, \textbf{Marrakech}, \textbf{Rabat}, \textbf{Tanger} et \textbf{Agadir}, avec une présence opérationnelle au niveau des aéroports (parmi lesquels Casablanca Mohammed~V, Marrakech Ménara et Agadir Al Massira).

\vspace{0.3cm}

\hspace{1.2em}Son fonctionnement repose sur un parcours simple et contrôlé : grâce à des \textbf{guichets intelligents} et des \textbf{bornes de prépaiement} installés dans les terminaux, le passager consulte un \textbf{tarif fixe} affiché à l’avance, règle la course avant de monter dans le véhicule, puis reçoit un ticket / reçu de prépaiement (souvent associé à un \textbf{QR code}). Ce mécanisme évite la négociation avec le chauffeur et réduit les litiges tarifaires. Le voyageur peut également réserver via le site \textbf{wetaxi.ma} ou l’application mobile.

\begin{figure}[H]
\centering
\safeimage{width=0.85\textwidth}{images/wetaxi_reception_aeroport.jpg}
\caption{Guichet / accueil WeTaxi dans un aéroport marocain}
\label{fig:wetaxi_reception}
\end{figure}

\hspace{1.2em}Selon les résultats communiqués par Tripass, WeTaxi a déjà transporté plus de \textbf{2,1~millions de passagers}, avec une note moyenne de satisfaction d’environ \textbf{4,3/5}. La plateforme intègre un système de \textbf{traçabilité numérique} des courses (suivi des trajets, évaluation des prestations) et a permis la restitution de \textbf{2\,341 objets} oubliés. Grâce au prépaiement systématique, \textbf{aucun litige tarifaire} n’aurait été enregistré depuis le lancement du service.

\vspace{0.3cm}

\hspace{1.2em}WeTaxi place également les \textbf{chauffeurs partenaires} au centre du dispositif : portefeuille numérique pour gérer les revenus, application dédiée au suivi des courses, et programme d’accompagnement à l’usage des outils numériques. Tripass vise principalement une clientèle composée de \textbf{touristes} et de voyageurs ayant besoin d’un transport fiable dès leur arrivée. L’entreprise compte environ \textbf{93 salariés}, et positionne WeTaxi comme la première brique d’un écosystème destiné à moderniser la chaîne de transport et d’accueil des voyageurs au Maroc.

\subsection{Organisation de l'entreprise}

L'organisation de Tripass est structurée autour de plusieurs directions complémentaires, comme illustré dans la Figure~\ref{fig:organigramme_Tripass}.

\begin{itemize}[leftmargin=2em,itemsep=4pt]
\item \textbf{Direction générale :} La direction générale de Tripass est assurée par \textbf{Abid Khirani}. Elle supervise l'ensemble des activités de l'entreprise et veille à la bonne coordination entre les différents départements.
\item \textbf{Direction développement \& technique :} La direction développement \& technique est assurée par \textbf{Amine Ahbouch}. Elle est chargée de la conception, de l'évolution et de la maintenance des solutions numériques proposées par Tripass, notamment la plateforme WeTaxi.
\item \textbf{Direction administrative :} La direction administrative est assurée par \textbf{Abderrahman Simou}. Elle s'occupe de la gestion administrative, de l'organisation interne et du suivi des aspects opérationnels de l'entreprise.
\end{itemize}

\vspace{0.3cm}

% Organigramme drawn with TikZ
\begin{figure}[H]
\centering
\safeimage{width=\textwidth,height=0.25\textheight}{images/organigramme de Tripass.png}\caption{L'organigramme de Tripass}
\label{fig:organigramme_Tripass}
\end{figure}


\section{Problématique et Objectifs du stage}

% ===================== 3.1 =====================
\subsection{Étude de l’existant}

Dans le domaine du marketing digital, plusieurs solutions professionnelles permettent de gérer les campagnes, les publications et les performances sur différents canaux numériques. L’étude de ces solutions permet d’identifier les fonctionnalités déjà disponibles sur le marché, mais aussi les limites qui justifient la conception d’une plateforme plus intégrée et intelligente.

L’analyse comparative montre que les solutions existantes couvrent efficacement certains besoins du marketing digital, notamment la planification de contenu, la gestion des réseaux sociaux, l’analyse des performances et l’assistance par intelligence artificielle. Toutefois, elles restent souvent spécialisées dans un domaine précis, limitées à certains écosystèmes ou conditionnées par des offres payantes pour les fonctionnalités avancées.

Par ailleurs, l’intégration avec différentes plateformes publicitaires et réseaux sociaux (Google Ads, TikTok Ads, LinkedIn, etc.) n’est pas toujours centralisée dans une seule interface. Les utilisateurs doivent ainsi naviguer entre plusieurs outils, ce qui peut entraîner une complexité accrue, une perte de temps et une difficulté à obtenir une vision globale et cohérente des performances. Ces limites mettent en évidence la nécessité de développer une solution intégrée, capable de centraliser les fonctionnalités, d’exploiter pleinement l’intelligence artificielle sur l’ensemble du processus marketing et d’offrir une expérience utilisateur plus fluide, cohérente et performante.

Le tableau~\ref{tab:solutions_existantes} ci-dessous synthétise cette comparaison des principales solutions du marché.

\begin{table}[hbp]
\centering
\caption{Comparaison des solutions existantes de marketing digital}
\label{tab:solutions_existantes}
\small
\renewcommand{\arraystretch}{1.2}
\begin{tabularx}{\textwidth}{|>{\bfseries\raggedright\arraybackslash}p{2.8cm}|>{\raggedright\arraybackslash}X|>{\raggedright\arraybackslash}X|}
\hline
\rowcolor{tableheadergreen}
\textcolor{white}{Solution} & \textcolor{white}{\textbf{Fonctionnalités principales}} & \textcolor{white}{\textbf{Limites}} \\
\hline
Hootsuite & Planification des publications, gestion de plusieurs comptes sociaux, analyse des performances et assistance IA pour la génération de contenu. & Solution principalement orientée réseaux sociaux ; certaines fonctionnalités avancées nécessitent des abonnements payants. \\
\hline
Buffer & Programmation de contenu, suivi de l’engagement, interface simple et assistance IA pour la rédaction et l’amélioration des publications. & Fonctionnalités analytiques et automatisations plus limitées pour une gestion marketing globale. \\
\hline
HubSpot & CRM, marketing automation, email marketing, gestion des campagnes, analyse des performances et outils IA pour le contenu et les données. & Plateforme complète mais complexe ; coût élevé et prise en main plus exigeante pour les petites structures. \\
\hline
Sprout Social & Planification avancée, analyse approfondie, collaboration en équipe et outils IA pour la création de contenu et l’analyse des interactions. & Outil puissant mais principalement centré sur les réseaux sociaux ; tarification élevée pour les fonctionnalités avancées. \\
\hline
\end{tabularx}
\end{table}

\FloatBarrier

Cette synthèse confirme l’intérêt de concevoir une plateforme unique, plus intégrée et intelligente, capable de répondre aux limites identifiées ci-dessus.


% ===================== 3.2 =====================
\subsection{Problématique}

L’essor du marketing digital et la multiplication des plateformes en ligne comme LinkedIn, Instagram, TikTok ou Google Ads ont profondément transformé la manière dont les entreprises interagissent avec leurs clients, tout en rendant la gestion des campagnes plus complexe et exigeante en temps, en compétences techniques et en analyse continue.

\vspace{0.3cm}

Malgré les avancées technologiques, des défis persistent, notamment la centralisation des données, l’automatisation des processus et l’optimisation des performances. Ces constats, observés durant mon stage, m’ont amené à formuler la problématique suivante : \textbf{comment concevoir une plateforme de gestion marketing multi-plateformes} permettant de :
\begin{itemize}[leftmargin=2em,itemsep=4pt]
\item intégrer un assistant conversationnel (agent IA) dédié aux actions marketing,
\item automatiser la création de contenu, la planification et les campagnes,
\item offrir une interface intuitive accessible à tous les niveaux techniques,
\item assurer une intégration fiable avec plusieurs plateformes,
\item garantir de bonnes performances et sécuriser les informations sensibles des utilisateurs et des campagnes ?
\end{itemize}

% ===================== 3.3 =====================
\subsection{Objectifs}

Le principal objectif de ce projet est de développer une plateforme web de gestion marketing multi-plateformes, nommée Agent-Market, capable de centraliser, automatiser et optimiser les stratégies de marketing digital tout en améliorant l’efficacité opérationnelle des entreprises. Plus précisément, les objectifs sont les suivants :

\begin{itemize}[leftmargin=2em,itemsep=6pt]

\item Créer une interface utilisateur intuitive : concevoir une interface ergonomique centralisant la gestion des campagnes, publications, performances et de l’assistant IA pour simplifier l’usage multi-plateformes.

\item Intégrer un assistant conversationnel (agent IA) : mettre à disposition un chat agent pour assister l’utilisateur dans ses activités marketing.

\item Automatiser les processus marketing : s’appuyer sur l’IA pour la génération de contenu, la planification des publications et l’orchestration de campagnes, afin de réduire les tâches manuelles.

\item Garantir la sécurité des données : protéger les informations utilisateurs, les tokens d’authentification (OAuth) et les données liées aux campagnes marketing.

\item Assurer l’intégration multi-plateformes et la maintenabilité : développer des connecteurs fiables vers les réseaux sociaux et outils publicitaires, au sein d’une architecture modulaire et évolutive.

\end{itemize}

% ===================== 3.4 =====================
\subsection{Motivations}

Le choix de ce projet s’explique par plusieurs motivations personnelles, professionnelles et organisationnelles :

\begin{itemize}[leftmargin=2em,itemsep=6pt]

\item Contexte du stage chez Tripass : l’entreprise évolue dans un environnement numérique où la communication et l’acquisition passent par les canaux digitaux, ce qui justifie une plateforme marketing centralisée enrichie d’un assistant IA.

\item Montée en compétences techniques : ce projet m’a permis d’approfondir le full-stack, l’intégration d’APIs, le traitement asynchrone et la mise en place d’un agent conversationnel.

\item Réponse à un besoin concret : plutôt que de naviguer entre plusieurs outils, l’utilisateur dispose d’un espace unique, complété par un chat agent IA, pour ses actions marketing.

\item Valeur ajoutée pour l’utilisateur final : offrir un assistant conversationnel intégré à la plateforme afin de gagner en efficacité et en clarté opérationnelle.

\end{itemize}

% ===============================================================
% SECTION 4 — GESTION DE PROJET
% ===============================================================

\section{Gestion de projet}

Pour le développement de la plateforme \textbf{Agent-Market} au sein de Tripass, j’ai opté pour la méthodologie \textbf{Waterfall}. Cette approche séquentielle m’a permis de réaliser chaque phase après la précédente, avec une documentation claire et un suivi rigoureux. Elle convenait bien au projet, car j’étais le principal responsable de la conception et du développement, sous l’encadrement de l’équipe Tripass.

% ===================== 4.1 =====================
\subsection{Méthodologie Waterfall}

La méthodologie de développement adoptée pour ce projet est la méthode Waterfall, une approche linéaire qui se déroule en plusieurs phases distinctes. La figure~\ref{fig:waterfall_method} en présente le principe général.

\begin{figure}[H]
\centering
\safeimage{width=1.0\textwidth,height=0.45\textheight}{images/watterfall.png}
\caption{Méthodologie Waterfall}
\label{fig:waterfall_method}
\end{figure}

Les phases mises en œuvre dans le cadre d’Agent-Market sont les suivantes :

\begin{itemize}[leftmargin=2em,itemsep=6pt]

\item \textbf{Requirements :} Identification et documentation des exigences spécifiques d’Agent-Market, incluant des échanges avec les parties prenantes de Tripass afin de définir clairement les besoins et attentes (gestion des publications, campagnes, leads et assistant IA). Cette phase critique assure une compréhension commune des objectifs du projet.

\item \textbf{Analysis :} Analyse détaillée des besoins et spécifications fonctionnelles, comprenant l’étude de faisabilité, l’évaluation des risques et la définition des critères de succès pour s’assurer que toutes les exigences sont réalisables. Cette étape permet de prévenir les problèmes potentiels avant qu’ils ne surviennent.

\item \textbf{Design :} Élaboration de l’architecture technique et des designs de l’application, incluant la création des schémas de bases de données, des diagrammes UML et des prototypes d’interface utilisateur, afin de fournir une feuille de route claire avant la phase de codage.

\item \textbf{Code :} Programmation et implémentation des fonctionnalités définies (frontend, API, worker, intégrations et assistant conversationnel), avec des revues régulières et l’utilisation de pratiques de développement standard pour garantir la qualité et la conformité aux spécifications.

\item \textbf{Test :} Vérification et validation des fonctionnalités développées par des tests unitaires, des tests d’intégration, des tests système et des tests d’acceptation, afin de s’assurer que la plateforme répond aux exigences initiales et reste fiable.

\item \textbf{Maintenance :} Suivi et correction des éventuels problèmes post-déploiement, incluant la gestion des incidents, les mises à jour et l’ajout de nouvelles fonctionnalités en réponse aux retours des utilisateurs. Cette phase assure la pérennité et l’amélioration continue de la solution.

\end{itemize}

% ===================== 1.4.2 =====================
\subsection{Diagramme de Gantt}

Le diagramme de Gantt, nommé d’après son créateur Henry Gantt, est une représentation visuelle des tâches et de leurs échéances. Il offre une vision claire du projet, facilitant l’organisation, la coordination et le suivi des activités.

Cet outil a joué un rôle important dans la progression du projet, en permettant de respecter les délais, de gérer les dépendances entre les tâches et de suivre l’avancement des travaux selon les phases Waterfall.

Dans le cadre d’Agent-Market, le diagramme de Gantt a servi à planifier le stage de quatre mois en découpant le travail par modules (recueil des besoins, conception UML, infrastructure, authentification et workspaces, publications sociales, campagnes publicitaires et leads, assistant IA, temps réel et reporting, tests, puis documentation et déploiement). Il a ainsi permis de visualiser les enchaînements entre ces lots, d’anticiper les chevauchements (par exemple entre le frontend et le worker) et d’assurer un suivi régulier de l’avancement auprès de l’équipe Tripass.

La figure~\ref{fig:gantt_project} ci-dessous illustre le diagramme de Gantt du projet :

\begin{figure}[H]
\centering
\safeimage{width=1.0\textwidth,height=0.25\textheight}{images/Gantt.png}
\caption{Diagramme de Gantt du projet}
\label{fig:gantt_project}
\end{figure}



% ===================== 4.4 =====================
\subsection{Conclusion}

En conclusion, ce chapitre a présenté l’organisme d’accueil Tripass, le contexte du projet Agent-Market, ainsi que la problématique, les objectifs et les motivations du stage. Il a également détaillé la méthode de gestion de projet retenue, à savoir la méthodologie Waterfall, ainsi que l’outil Notion utilisé pour le suivi et le diagramme de Gantt pour la planification temporelle.

Cette vision d’ensemble permet de mieux situer le projet dans son environnement organisationnel et méthodologique. Les chapitres suivants approfondiront l’analyse du projet, la conception et la réalisation de la plateforme.



\newpage
% ===============================================================
% PAGE 22 — CHAPTER 2
% ===============================================================

\refstepcounter{chapter}

\addcontentsline{toc}{chapter}{Chapitre \thechapter : Analyse du projet}

\chapterpagetitle{Analyse du projet}

% ===============================================================
% PAGE 23
% ===============================================================


\section{Introduction}

Dans ce chapitre, nous analysons le projet \textbf{Agent-Market} à partir de la problématique et des objectifs définis au chapitre~1. Nous y formalisons les besoins fonctionnels, non fonctionnels et techniques, puis présentons les fonctionnalités de l’assistant conversationnel (agent IA).



\section{Spécification des besoins}

\subsection{Besoins fonctionnels}

\textbf{Gestion des workspaces et authentification :}

\begin{itemize}[leftmargin=2em,itemsep=6pt]
\item[] \textbf{--- Authentification sécurisée :} L’application doit permettre aux utilisateurs de s’authentifier de manière fiable et de gérer leurs sessions au sein d’un espace de travail partagé.
\begin{itemize}[label=$\checkmark$,leftmargin=2em,itemsep=3pt]
\item \textbf{Inscription et connexion :} Créer un compte, se connecter et se déconnecter via JWT.
\item \textbf{Gestion des workspaces :} Isoler les données marketing par organisation (membres, rôles, paramètres).
\item \textbf{Protection des routes :} Restreindre l’accès aux ressources selon le rôle et le workspace actif.
\end{itemize}

\item[] \textbf{--- Gestion d’équipe :} Permettre la collaboration entre plusieurs utilisateurs au sein d’un même workspace.
\begin{itemize}[label=$\checkmark$,leftmargin=2em,itemsep=3pt]
\item \textbf{Invitation de membres :} Ajouter des collaborateurs avec des droits adaptés.
\item \textbf{Profil utilisateur :} Modifier les informations personnelles et les préférences.
\end{itemize}
\end{itemize}

\textbf{Gestion sociale et publication multi-canal :}

\begin{itemize}[leftmargin=2em,itemsep=6pt]
\item[] \textbf{--- Connexion des comptes sociaux :} Établir une connexion entre Agent-Market et les plateformes sociales pour une publication centralisée.
\begin{itemize}[label=$\checkmark$,leftmargin=2em,itemsep=3pt]
\item \textbf{OAuth :} Authentifier de manière sécurisée LinkedIn, Meta (Facebook/Instagram), TikTok et X (Twitter).
\item \textbf{Stockage des tokens :} Conserver et rafraîchir automatiquement les jetons d’accès.
\item \textbf{Gestion multi-comptes :} Associer plusieurs comptes sociaux à un même workspace.
\end{itemize}

\item[] \textbf{--- Création et édition de publications :} Permettre la rédaction de contenus marketing adaptés à chaque canal.
\begin{itemize}[label=$\checkmark$,leftmargin=2em,itemsep=3pt]
\item \textbf{Brouillons :} Enregistrer des posts avant publication.
\item \textbf{Médias :} Joindre images et vidéos via la bibliothèque média.
\item \textbf{Prévisualisation :} Visualiser le rendu par plateforme avant diffusion.
\end{itemize}

\item[] \textbf{--- Génération de contenu par IA :} Automatiser la production de textes marketing optimisés.
\begin{itemize}[label=$\checkmark$,leftmargin=2em,itemsep=3pt]
\item \textbf{Adaptation par réseau :} Ajuster le ton et le format selon LinkedIn, Instagram, TikTok, etc.
\item \textbf{Suggestions intelligentes :} Proposer hashtags, accroches et variantes de contenu.
\item \textbf{Catalogue IA :} Configurer les fournisseurs et modèles LLM par workspace.
\end{itemize}

\item[] \textbf{--- Planification et publication :} Programmer et exécuter les publications de manière fiable.
\begin{itemize}[label=$\checkmark$,leftmargin=2em,itemsep=3pt]
\item \textbf{Calendrier éditorial :} Définir date et heure de publication.
\item \textbf{Publication multi-canal :} Diffuser simultanément sur plusieurs comptes.
\item \textbf{Automatisation :} Exécuter les publications planifiées via Celery Beat et le worker Python.
\end{itemize}

\item[] \textbf{--- Suivi des performances sociales :} Importer et afficher les métriques d’engagement.
\begin{itemize}[label=$\checkmark$,leftmargin=2em,itemsep=3pt]
\item \textbf{Import des followers et posts :} Synchroniser les données depuis les API externes.
\item \textbf{Indicateurs :} Suivre likes, commentaires, partages et impressions.
\item \textbf{Rapports :} Consolider les performances dans le dashboard.
\end{itemize}
\end{itemize}

\textbf{Campagnes publicitaires Google Ads et Meta Ads :}

\begin{itemize}[leftmargin=2em,itemsep=6pt]
\item[] \textbf{--- Connexion des comptes publicitaires :} Relier les comptes Google Ads et Meta Ads au workspace.
\begin{itemize}[label=$\checkmark$,leftmargin=2em,itemsep=3pt]
\item \textbf{OAuth Google / Meta :} Authentifier les comptes annonceurs de manière sécurisée.
\item \textbf{Sélection des comptes :} Choisir les comptes et pages à piloter depuis l’interface.
\end{itemize}

\item[] \textbf{--- Création et gestion de campagnes :} Permettre la configuration complète des campagnes publicitaires.
\begin{itemize}[label=$\checkmark$,leftmargin=2em,itemsep=3pt]
\item \textbf{Campagnes Google Ads :} Créer campagnes, groupes d’annonces, mots-clés et annonces RSA.
\item \textbf{Campagnes Meta Ads :} Configurer campagnes, ensembles de publicités, ciblage et créatifs.
\item \textbf{Budget et planning :} Définir budgets quotidiens, dates de début et de fin.
\end{itemize}

\item[] \textbf{--- Synchronisation et suivi :} Garantir une vision cohérente des campagnes distantes et locales.
\begin{itemize}[label=$\checkmark$,leftmargin=2em,itemsep=3pt]
\item \textbf{Synchronisation automatique :} Importer périodiquement l’état des campagnes via Celery.
\item \textbf{Insights :} Afficher CTR, CPC, conversions, dépenses et autres KPIs.
\item \textbf{Workflow guidé :} Guider l’utilisateur étape par étape dans les assistants de création.
\end{itemize}
\end{itemize}

\textbf{Agent IA d’orchestration marketing :}

\begin{itemize}[leftmargin=2em,itemsep=6pt]
\item[] \textbf{--- Assistant conversationnel :} Le chatbot doit comprendre les intentions marketing et proposer des actions concrètes.
\begin{itemize}[label=$\checkmark$,leftmargin=2em,itemsep=3pt]
\item \textbf{Compréhension de la requête :} Analyser les demandes en langage naturel (création de campagne, publication, analyse).
\item \textbf{Planification de workflow :} Décomposer une tâche complexe en étapes exécutables.
\item \textbf{Historique :} Conserver le contexte de la conversation pour des réponses cohérentes.
\end{itemize}

\item[] \textbf{--- Exécution des workflows :} L’agent doit pouvoir déclencher et suivre des actions sur la plateforme.
\begin{itemize}[label=$\checkmark$,leftmargin=2em,itemsep=3pt]
\item \textbf{Orchestration :} Enchaîner appels API, tâches Celery et outils publicitaires.
\item \textbf{Suivi en temps réel :} Afficher l’avancement de chaque étape du workflow.
\item \textbf{Résultats structurés :} Présenter les sorties (créatifs, métriques, erreurs) de manière lisible.
\end{itemize}

\item[] \textbf{--- Gestion des demandes multiples :} Le système doit traiter plusieurs workflows sans dégradation notable.
\begin{itemize}[label=$\checkmark$,leftmargin=2em,itemsep=3pt]
\item \textbf{Files d’attente :} Utiliser Redis et Celery pour absorber la charge.
\item \textbf{Annulation :} Permettre l’arrêt d’un workflow en cours si nécessaire.
\end{itemize}
\end{itemize}

\textbf{Landing pages et génération de leads :}

\begin{itemize}[leftmargin=2em,itemsep=6pt]
\item[] \textbf{--- Création de landing pages :} Faciliter la mise en ligne de pages de conversion liées aux campagnes.
\begin{itemize}[label=$\checkmark$,leftmargin=2em,itemsep=3pt]
\item \textbf{Éditeur :} Créer et personnaliser des pages avec slug public.
\item \textbf{Formulaires :} Collecter nom, email et champs personnalisés.
\item \textbf{Tracking UTM :} Associer les leads aux campagnes sources.
\end{itemize}

\item[] \textbf{--- Gestion des leads :} Centraliser et exploiter les contacts générés.
\begin{itemize}[label=$\checkmark$,leftmargin=2em,itemsep=3pt]
\item \textbf{Liste et filtres :} Consulter, rechercher et exporter les leads.
\item \textbf{Notifications :} Alerter l’équipe lors de nouvelles soumissions.
\end{itemize}
\end{itemize}

\textbf{Dashboard, reporting et suivi budgétaire :}

\begin{itemize}[leftmargin=2em,itemsep=6pt]
\item[] \textbf{--- Tableau de bord unifié :} Offrir une vue consolidée de l’activité marketing.
\begin{itemize}[label=$\checkmark$,leftmargin=2em,itemsep=3pt]
\item \textbf{KPIs globaux :} Afficher dépenses, performances sociales et publicitaires.
\item \textbf{Graphiques :} Visualiser l’évolution des indicateurs dans le temps.
\item \textbf{Suivi budgétaire :} Comparer budgets alloués et dépenses réelles.
\end{itemize}
\end{itemize}

\textbf{Temps réel et notifications :}

\begin{itemize}[leftmargin=2em,itemsep=6pt]
\item[] \textbf{--- Diffusion d’événements live :} Informer l’interface des changements dès qu’ils surviennent.
\begin{itemize}[label=$\checkmark$,leftmargin=2em,itemsep=3pt]
\item \textbf{WebSocket :} Transmettre les événements via SocketCluster.
\item \textbf{Redis Pub/Sub :} Relier l’API, le worker et le service temps réel.
\item \textbf{Fil d’activité :} Afficher publications, campagnes et alertes en direct sur le dashboard.
\end{itemize}
\end{itemize}

\subsection{Besoins non fonctionnels}

\begin{itemize}[leftmargin=2em,itemsep=8pt]

\item[$\checkmark$] \textbf{Performance :} \\
Assurer des temps de réponse rapides pour la consultation du dashboard, la génération de contenu IA, la synchronisation des campagnes publicitaires et l’exécution des publications planifiées.

\item[$\checkmark$] \textbf{Sécurité :} \\
Implémenter des mécanismes d’authentification et d’autorisation robustes (JWT, OAuth), le chiffrement des configurations sensibles et la protection des routes internes worker (secret partagé).

\item[$\checkmark$] \textbf{Fiabilité :} \\
Garantir un traitement fiable des tâches asynchrones, des sauvegardes régulières des données MongoDB et une haute disponibilité des services critiques (API, worker, temps réel).

\item[$\checkmark$] \textbf{Maintenabilité :} \\
Développer une architecture modulaire (API Fastify, worker Celery, service realtime) et une documentation technique claire pour faciliter les évolutions futures.

\item[$\checkmark$] \textbf{Évolutivité :} \\
Concevoir le système pour qu’il puisse évoluer avec l’augmentation du nombre d’utilisateurs, de campagnes, de publications et d’interactions avec l’agent IA, sans compromettre les performances.

\item[$\checkmark$] \textbf{Convivialité :} \\
Assurer une interface utilisateur intuitive centralisant publications, campagnes, leads et assistant IA, afin de réduire la complexité multi-plateformes.

\item[$\checkmark$] \textbf{Intégrité des données :} \\
Garantir la cohérence entre les données locales (MongoDB) et les plateformes externes lors des synchronisations publicitaires et sociales.

\end{itemize}

\subsection{Besoins techniques}

Les besoins techniques ci-dessous résument les composants indispensables pour satisfaire les exigences ci-dessus. Leur rôle détaillé et leur justification seront approfondis au chapitre consacré aux outils et technologies de développement.

\begin{table}[H]
\centering
\caption{Synthèse des besoins techniques}
\label{tab:besoins_techniques}
\renewcommand{\arraystretch}{1.4}
\begin{tabular}{|p{0.22\textwidth}|p{0.68\textwidth}|}
\hline
\textbf{Couche} & \textbf{Technologies requises} \\
\hline
API & Fastify, Node.js, MongoDB, JWT, ioredis (cache, pub/sub, pont Celery) \\
\hline
Worker & Python 3.11, Celery, Celery Beat, connecteurs Google Ads / Meta / réseaux sociaux, LLM (Gemini, OpenAI, Anthropic) \\
\hline
Temps réel & SocketCluster, Redis Pub/Sub \\
\hline
Frontend & Next.js (App Router), React, CSS global, client WebSocket \\
\hline
Infrastructure & MongoDB~7, Redis~7, Docker Compose \\
\hline
\end{tabular}
\end{table}

\section{Présentation des fonctionnalités de l’assistant IA}

Dans le cadre du développement d’\textbf{Agent-Market}, un assistant conversationnel (agent IA) a été conçu pour accompagner les utilisateurs dans leurs activités marketing : création de contenu, planification de publications, orchestration de campagnes et suivi des performances. Plusieurs fonctionnalités ont été mises en place pour améliorer l’efficacité opérationnelle et offrir une expérience plus fluide.

Pour les utilisateurs :

\begin{itemize}[leftmargin=2em,itemsep=6pt]

\item \textbf{Conversation en langage naturel :} L’utilisateur exprime son besoin marketing dans un chat dédié (par exemple créer une publication, préparer une campagne ou analyser des performances). Cette entrée constitue le point de départ de l’assistance IA.

\item \textbf{Contexte workspace et historique :} L’assistant s’appuie sur le workspace actif, les comptes connectés et l’historique de conversation afin de produire des réponses cohérentes et adaptées à l’environnement de l’utilisateur.

\item \textbf{Planification de workflows :} À partir de la requête, l’agent décompose la demande en étapes exécutables (génération de contenu, sélection de comptes, planification, publication, etc.) et propose un plan d’action structuré.

\item \textbf{Validation avant exécution :} Le workflow proposé est présenté à l’utilisateur, qui peut le valider, l’ajuster ou l’annuler. Aucune action marketing sensible n’est lancée sans cette validation.

\item \textbf{Exécution assistée :} Une fois validé, le workflow est pris en charge par la plateforme (API, worker Celery et connecteurs vers les réseaux sociaux ou outils publicitaires) pour réaliser les actions prévues.

\item \textbf{Suivi et restitution des résultats :} L’utilisateur peut suivre l’avancement des étapes et consulter les résultats obtenus (contenus générés, publications, campagnes préparées ou erreurs éventuelles) directement depuis l’interface de l’assistant.

\end{itemize}

Une attention particulière a également été portée à l’intégration multi-fournisseurs IA (OpenAI, Gemini, Anthropic), à la fluidité de l’expérience conversationnelle et à la cohérence des données manipulées par l’agent au sein du workspace.

\section{Conclusion}

Ce chapitre a analysé le projet \textbf{Agent-Market} à travers ses besoins fonctionnels, non fonctionnels et techniques, ainsi que les fonctionnalités de l’assistant IA. Le chapitre suivant traduira ces éléments en modèles UML (cas d’utilisation, séquences, classes, activités) ; le chapitre sur les outils et technologies détaillera ensuite le rôle de chaque composant de la stack.


\newpage

% ===============================================================
% CHAPTER 3
% ===============================================================

\refstepcounter{chapter}

\addcontentsline{toc}{chapter}{Chapitre \thechapter : Étude conceptuelle}

\chapterpagetitle{Étude conceptuelle}

% Espacement flottants : éviter les grands vides autour des figures
\setlength{\intextsep}{8pt plus 2pt minus 2pt}
\setlength{\textfloatsep}{10pt plus 2pt minus 2pt}
\setlength{\floatsep}{8pt plus 2pt minus 2pt}
\captionsetup{skip=6pt}

% Garde la phrase d'accroche avec la figure (évite une page qui commence par une figure)
% Usage : \figlead[0.55\textheight]{phrase...}
\newcommand{\figlead}[2][0.82\textheight]{%
  \par\Needspace{#1}%
  #2\par\nopagebreak\vspace{2pt}%
}
% Colle le titre « Description textuelle » à son tableau (sans Needspace agressif)
\newcommand{\deschead}[1]{%
  \par\noindent\textbf{Description textuelle «~#1~»}\par\nopagebreak\vspace{2pt}%
}

% ===============================================================
% CHAPITRE — MODÉLISATION UML
% ===============================================================

\section{Introduction}

L’analyse et la conception du projet \textbf{Agent-Market} ont été menées dans le but de définir de manière précise les besoins fonctionnels et techniques de la solution. Cette phase pose les bases du développement en s’appuyant sur la modélisation UML : cas d’utilisation, séquences, classes et activités.

\begin{figure}[H]
\centering
\safeimage{width=0.45\textwidth, height=0.18\textheight}{images/uml.png}
\end{figure}

Les diagrammes présentés dans ce chapitre décrivent à la fois la structure du système et son comportement dynamique, en cohérence avec les modules métier de la plateforme.

\section{Le langage UML}

UML fournit un ensemble de notations graphiques et de conventions permettant de modéliser les aspects structurels et comportementaux d’un système. Il facilite la compréhension globale du système en offrant une représentation visuelle des composants, de leurs relations et de leurs interactions.

De plus, UML améliore la communication entre les différents acteurs du projet (développeurs, chefs de projet, clients) et constitue un support essentiel pour la conception et la documentation des systèmes logiciels complexes.

\section{Choix du langage UML}

Le choix d’UML comme langage de modélisation est justifié par plusieurs avantages :

\begin{itemize}[leftmargin=2em,itemsep=6pt]

\item \textbf{Standardisation :} UML est un standard international largement adopté, garantissant une uniformité dans la modélisation des systèmes.

\item \textbf{Compréhension visuelle :} Les diagrammes UML offrent une représentation graphique claire facilitant la compréhension des concepts et des relations.

\item \textbf{Documentation :} UML permet de documenter efficacement les différentes phases du développement logiciel.

\item \textbf{Flexibilité :} Il propose plusieurs types de diagrammes adaptés aux différents besoins (structure, comportement, interactions).

\item \textbf{Collaboration :} Il favorise une meilleure communication et coordination entre les membres de l’équipe projet.

\end{itemize}

% ===============================================================
\section{Diagramme de cas d’utilisation}

Le diagramme de cas d’utilisation représente les interactions entre les acteurs et le système. Il met en évidence les fonctionnalités principales d’Agent-Market et les scénarios d’usage associés aux modules métier.

\subsection{Identification des acteurs}

Les acteurs se répartissent entre \textbf{humains} (Utilisateur, Admin, Visiteur web) et \textbf{systèmes externes} (Fournisseur IA, plateformes sociales et publicitaires). Les premiers initient les cas d’usage ; les seconds réalisent les échanges techniques (IA, OAuth, publication). Le tableau~\ref{tab:acteurs} en synthétise les rôles.

\begin{table}[H]
\centering
\caption{Identification des acteurs du système}
\label{tab:acteurs}
\renewcommand{\arraystretch}{1.35}
\begin{tabularx}{\textwidth}{|>{\bfseries\raggedright\arraybackslash}p{4.5cm}|>{\raggedright\arraybackslash}X|}
\hline
\rowcolor{tableheadergreen}
\textcolor{white}{Acteur} & \textcolor{white}{\textbf{Rôle}} \\
\hline
Utilisateur & Marketeur authentifié : publications, campagnes, profil et dialogue avec l’assistant IA. \\
\hline
Admin & Spécialisation de l’Utilisateur ; gère le workspace (membres, rôles, préférences). \\
\hline
Fournisseur IA & Service externe (ex.\ Gemini) pour comprendre les demandes, planifier et exécuter. \\
\hline
Plateformes sociales & Réseaux (Meta, X, etc.) pour OAuth, publication et métriques. \\
\hline
Plateformes publicitaires & Google Ads / Meta Ads pour OAuth, publication et performances. \\
\hline
Visiteur web & Internautes soumettant un formulaire public de lead. \\
\hline
\end{tabularx}
\end{table}

Cette distinction permet de clarifier les frontières du système : les acteurs humains initient les cas d’usage, tandis que les systèmes externes réalisent les échanges techniques indispensables à l’automatisation marketing.

\subsection{Diagramme global des cas d’utilisation}

La figure~\ref{fig:uc_global} présente la vue d’ensemble d’Agent-Market. L’\textbf{Utilisateur} accède aux modules métier (assistant IA, publications, campagnes, profil) ; l’\textbf{Admin} hérite de ces droits et ajoute la \textbf{Gestion Workspace}. Tous ces cas \textbf{incluent} «~Se connecter~». Les \textbf{extend} de la gestion workspace couvrent invitation, rôles et préférences.

\begin{figure}[H]
\centering
\safeimage{width=0.92\textwidth,height=0.44\textheight}{images/use_case_global.png}
\caption{Diagramme global des cas d’utilisation du système Agent-Market}
\label{fig:uc_global}
\end{figure}

Le tableau détaille le cas «~Authentification et gestion du workspace~» : connexion JWT, rattachement à un workspace, et actions Admin (membres, rôles, préférences). 

\deschead{Authentification et gestion du workspace}
\begin{table}[H]
\centering
\caption{Description textuelle — Authentification et Workspace}
\renewcommand{\arraystretch}{1.4}
\begin{tabular}{|p{0.28\textwidth}|p{0.62\textwidth}|}
\hline
\textbf{Cas d’utilisation} & Authentification et gestion du workspace \\
\hline
\textbf{Acteurs} & Utilisateur, Admin \\
\hline
\textbf{Résumé} &
Créer un compte, se connecter, rejoindre ou administrer un workspace (membres, rôles, préférences). \\
\hline
\textbf{Précondition} &
Accès à la plateforme. \\
\hline
\textbf{Scénario principal} &
\begin{itemize}[leftmargin=*,itemsep=2pt,topsep=2pt]
\item L’utilisateur se connecte (cas inclus par les autres modules).
\item Il crée un compte ou ouvre une session JWT.
\item L’Admin invite des membres, gère rôles et préférences.
\end{itemize}
\\

\hline
\end{tabular}
\end{table}

Les alternatives couvrent identifiants invalides, invitations expirées et permissions insuffisantes.

\subsection{Cas d’utilisation — Gestion des publications}

\figlead[0.55\textheight]{La figure~\ref{fig:uc_post} reprend le cycle brouillon–programmation–publication multi-comptes, ainsi que le recours optionnel à la génération de contenu par IA}
\begin{figure}[H]
\centering
\safeimage{width=1.0\textwidth,height=0.45\textheight}{images/use_case_post.png}
\caption{Diagramme de cas d’utilisation — Gestion des publications}
\label{fig:uc_post}
\end{figure}


\deschead{Gestion des publications}
\begin{table}[H]
\centering
\caption{Description textuelle — Gestion des publications}
\renewcommand{\arraystretch}{1.4}
\begin{tabular}{|p{0.28\textwidth}|p{0.62\textwidth}|}
\hline
\textbf{Cas d’utilisation} & Gestion des publications \\
\hline
\textbf{Acteurs} & Utilisateur, Fournisseur IA, Plateformes sociales \\
\hline
\textbf{Résumé} &
Créer, générer (IA), programmer et publier des posts multi-comptes, puis consulter calendrier et métriques. \\
\hline
\textbf{Précondition} &
Utilisateur authentifié. \\
\hline
\textbf{Scénario principal} &
\begin{itemize}[leftmargin=*,itemsep=2pt,topsep=2pt]
\item Connexion des comptes sociaux via OAuth.
\item Création d’un brouillon (éventuellement généré par IA).
\item Publication immédiate ou programmation.
\item Consultation du calendrier et des métriques.
\end{itemize}
\\
\hline
\end{tabular}
\end{table}

Le tableau détaille le cas «~Gestion des publications~» qui couvre génération IA, brouillon, programmation, publication, calendrier et métriques. Programmer et publier \textbf{incluent} le brouillon et la sélection de comptes ; cette dernière \textbf{inclut} l’OAuth. Les acteurs externes sont le \textbf{Fournisseur IA} et les \textbf{plateformes de médias sociaux}.

\subsection{Cas d’utilisation — Gestion des campagnes publicitaires}

La figure~\ref{fig:uc_ad} présente ces cas et relations pour la gestion des campagnes publicitaires~: génération IA, programmation et publication de campagnes, connexion OAuth Ads, performances, landing pages et leads. Programmer ou publier \textbf{inclut} le brouillon et la sélection d’un compte Ads. Le \textbf{Visiteur web} soumet le formulaire public de lead. 

\begin{figure}[H]
\centering
\safeimage{width=1.1\textwidth,height=0.6\textheight}{images/use_case_ad.png}
\caption{Diagramme de cas d’utilisation — Gestion des campagnes publicitaires}
\label{fig:uc_ad}
\end{figure}

On y retrouve la chaîne brouillon–programmation–publication, la connexion aux plateformes Ads et le parcours lead via landing page.

Le tableau détaille le cas «~Gestion des campagnes publicitaires~» : création, programmation et publication, gestion des landing pages et leads, analyse des performances.

\deschead{Gestion des campagnes publicitaires}
\begin{table}[H]
\centering
\caption{Description textuelle — Gestion des campagnes publicitaires}
\renewcommand{\arraystretch}{1.4}
\begin{tabular}{|p{0.28\textwidth}|p{0.62\textwidth}|}
\hline
\textbf{Cas d’utilisation} & Gestion des campagnes publicitaires \\
\hline
\textbf{Acteurs} & Utilisateur, Fournisseur IA, Plateformes Ads, Visiteur web \\
\hline
\textbf{Résumé} &
Créer, programmer et publier des campagnes, gérer landing pages et leads, analyser les performances. \\
\hline
\textbf{Précondition} &
Au moins un compte publicitaire connecté (ou connexion OAuth). \\
\hline
\textbf{Scénario principal} &
\begin{itemize}[leftmargin=*,itemsep=2pt,topsep=2pt]
\item Création d’une campagne en brouillon (contenu éventuellement IA).
\item Sélection du compte Ads, puis publication ou programmation.
\item Gestion des landing pages ; le visiteur soumet un lead.
\item Consultation des leads et des performances.
\end{itemize}
\\
\hline
\end{tabular}
\end{table}

Les alternatives couvrent l’échec OAuth Ads, le rejet de campagne et le formulaire de lead invalide.

\subsection{Cas d’utilisation de l’Assistant IA}

\figlead[0.40\textheight]{La figure~\ref{fig:uc_assistant} synthétise ce parcours conversationnel et décisionnel de l’Assistant IA.}
\begin{figure}[H]
\centering
\safeimage{width=0.75\textwidth,height=0.3\textheight}{images/use_case_assistant.png}
\caption{Diagramme de cas d’utilisation de l’Assistant IA}
\label{fig:uc_assistant}
\end{figure}

L’utilisateur \textbf{discute} avec l’assistant et \textbf{approuve} les workflows. Le fournisseur IA participe à la compréhension (toute langue), à la planification multi-étapes et à l’exécution. Les relations \textbf{include} enchaînent discussion $\rightarrow$ planification $\rightarrow$ approbation $\rightarrow$ exécution.


\deschead{Assistant IA}
\begin{table}[H]
\centering
\caption{Description textuelle — Assistant IA}
\renewcommand{\arraystretch}{1.4}
\begin{tabular}{|p{0.28\textwidth}|p{0.62\textwidth}|}
\hline
\textbf{Cas d’utilisation} & Assistant IA \\
\hline
\textbf{Acteurs} & Utilisateur, Fournisseur IA \\
\hline
\textbf{Résumé} &
Dialoguer en langage naturel, obtenir un plan multi-étapes, l’approuver puis l’exécuter via l’IA. \\
\hline
\textbf{Précondition} &
Utilisateur authentifié avec accès à l’assistant. \\
\hline
\textbf{Scénario principal} &
\begin{itemize}[leftmargin=*,itemsep=2pt,topsep=2pt]
\item L’utilisateur ouvre une conversation et exprime son besoin.
\item Le fournisseur IA comprend la demande et propose un plan.
\item L’utilisateur approuve le workflow.
\item Le système exécute les étapes et restitue le résultat.
\end{itemize}
\\
\hline
\textbf{Scénarios alternatifs} &
Demande ambiguë ; refus du plan ; échec d’une étape ; IA indisponible. \\
\hline
\end{tabular}
\end{table}



Ce module relie ainsi l’intention exprimée en langage naturel à une exécution outillée, sous contrôle explicite de l’utilisateur.

% ===============================================================
% 5. DIAGRAMMES DE SÉQUENCE
% ===============================================================

\section{Diagrammes de séquence}

Les diagrammes de séquence décrivent l’enchaînement chronologique des interactions entre acteurs et composants. Ils précisent le comportement dynamique des flux d’authentification, d’OAuth, de publication sociale et d’assistance IA. Contrairement aux cas d’utilisation, ils rendent visibles les messages échangés, les retours et les alternatives (\texttt{alt}, \texttt{loop}) entre l’interface, l’API, la base de données, le worker et les services externes.

\subsection{Authentification et contexte workspace}

Le parcours d’accès se découpe en trois phases : \textbf{inscription} (unicité de l’email, puis invitation ou workspace personnel), \textbf{connexion} (vérification du compte et chargement des workspaces) et \textbf{consultation} dans le workspace actif (contrôle de session puis lecture des données). Les participants sont l’utilisateur, le frontend Next.js, l’API backend et la base de données.

\figlead[0.88\textheight]{La figure~\ref{fig:seq_auth} présente ce déroulement sous forme de diagramme de séquence.}
\begin{figure}[H]
\centering
\safeimage{width=1.0\textwidth,height=0.85\textheight}{images/Authentification_et_contexte_workspace.png}
\caption{Diagramme de séquence — Authentification et contexte workspace}
\label{fig:seq_auth}
\end{figure}

Lors de l’inscription, le système vérifie que l’email n’existe pas déjà, crée le compte, puis soit valide une invitation et rattache l’utilisateur à un workspace existant, soit crée un workspace personnel. La connexion recharge les workspaces accessibles et fixe le workspace actif dans la session. Toute consultation ultérieure commence par un contrôle de session et de contexte, afin d’isoler les données multi-tenant avant la lecture en base.

En synthèse, le diagramme met en évidence :
\begin{itemize}[leftmargin=2em,itemsep=2pt,topsep=4pt]
\item la création de session JWT et le choix du workspace actif ;
\item le branchement invitation / workspace personnel à l’inscription ;
\item le garde-fou systématique avant toute lecture métier.
\end{itemize}

Ce scénario garantit qu’aucune consultation métier n’a lieu sans session valide et sans contexte workspace résolu. Il constitue le prérequis de tous les enchaînements suivants : sans workspace actif, ni OAuth ni publication ne peuvent être exécutés de manière cohérente.

\subsection{Connexion d’un compte social (OAuth)}

La connexion OAuth enchaîne quatre phases : \textbf{démarrage}, \textbf{autorisation} auprès de la plateforme externe, \textbf{enregistrement} du compte (avec sélection éventuelle de page/compte), puis \textbf{synchronisation} asynchrone (import posts et métriques par le worker). Les acteurs externes typiques sont Meta, X ou LinkedIn.

Au démarrage, l’interface demande à l’API une URL d’autorisation ; le contexte (workspace, réseau choisi) est mémorisé côté serveur. Après authentification chez le fournisseur, le callback renvoie un code ou des jetons. Lorsque plusieurs pages ou comptes sont disponibles, l’utilisateur en sélectionne un avant l’enregistrement en base. Le worker lance ensuite un import initial, persiste les données récupérées et notifie l’interface de la fin de synchronisation, ce qui actualise la liste des comptes connectés.

Les points clés à retenir sont :
\begin{itemize}[leftmargin=2em,itemsep=2pt,topsep=4pt]
\item la redirection OAuth et la mémorisation du contexte applicatif ;
\item la sélection optionnelle d’une page ou d’un compte côté fournisseur ;
\item l’import asynchrone (posts, métriques) avant la notification UI.
\end{itemize}

\figlead[0.88\textheight]{La figure~\ref{fig:seq_oauth} détaille ces interactions jusqu’à la notification de fin de synchronisation.}
\begin{figure}[H]
\centering
\safeimage{width=1.0\textwidth,height=0.8\textheight}{images/Connecter_compte.png}
\caption{Diagramme de séquence — Connexion d’un compte social (OAuth)}
\label{fig:seq_oauth}
\end{figure}

À l’issue de la synchronisation, l’API notifie l’interface afin d’actualiser la liste des comptes connectés. Cette étape conditionne la sélection multi-comptes lors de la création et de la publication d’un post.

\subsection{Planification et publication d’un post social}

\figlead[0.88\textheight]{La figure~\ref{fig:seq_post} formalise cette délégation au worker et le suivi du statut par compte.}
\begin{figure}[H]
\centering
\safeimage{width=1.0\textwidth,height=0.83\textheight}{images/Planification_et_publication_post_social.png}
\caption{Diagramme de séquence — Planification et publication d’un post social}
\label{fig:seq_post}
\end{figure}

Ce flux couvre le cycle \textbf{brouillon $\rightarrow$ immédiat ou programmé $\rightarrow$ multi-comptes}. Après création du brouillon, l’utilisateur choisit une publication immédiate (worker lancé tout de suite) ou programmée (planification persistée, exécution à l’heure due). L’exécution charge le post, publie en boucle sur chaque compte cible, enregistre les statuts, puis notifie l’interface.

La phase de création persiste le contenu, les médias et les comptes cibles comme brouillon. En publication immédiate, l’API délègue au worker sans attendre ; en publication programmée, seule la date/heure est enregistrée jusqu’à l’échéance. Pendant l’exécution, chaque compte fait l’objet d’un appel distinct vers le réseau social : succès et erreurs sont historisés séparément, puis un statut global est calculé. La notification finale permet à l’utilisateur de consulter le résultat sans recharger manuellement l’ensemble du dashboard.

Le diagramme insiste donc sur :
\begin{itemize}[leftmargin=2em,itemsep=2pt,topsep=4pt]
\item l’alternative publication immédiate / publication programmée ;
\item la boucle multi-comptes avec statut unitaire puis statut global ;
\item la remontée du résultat vers l’interface (notification).
\end{itemize}



Ce diagramme est indispensable pour comprendre une publication multi-réseaux fiable et traçable.

\subsection{Fonctionnement de l’Assistant IA}

Le dialogue suit la chaîne utilisateur $\rightarrow$ interface $\rightarrow$ API $\rightarrow$ worker IA $\rightarrow$ LLM. Après compréhension de la demande, une alternative distingue la \textbf{conversation simple} (réponse affichée) du \textbf{workflow} (plan à valider, exécution bouclée sur les plateformes externes, résumé LLM, résultat final).

Dans le cas conversationnel, le worker retourne directement une réponse exploitable dans le chat. Dans le cas workflow, un plan multi-étapes est soumis à validation : l’utilisateur approuve, l’API lance l’exécution, chaque étape interroge les connecteurs ou le LLM, puis un résumé consolide les résultats. Cette séparation protège les actions à effet de bord (publication, campagne) derrière un contrôle humain explicite.

On retient notamment :
\begin{itemize}[leftmargin=2em,itemsep=2pt,topsep=4pt]
\item la bifurcation conversation simple / workflow après analyse LLM ;
\item l’étape d’approbation avant toute exécution outillée ;
\item la génération d’un résumé final à l’issue de la boucle d’étapes.
\end{itemize}

\figlead[0.85\textheight]{La figure~\ref{fig:seq_assistant} montre ces deux branches au sein d’un même enchaînement.}
\begin{figure}[H]
\centering
\safeimage{width=0.9\textwidth,height=0.75\textheight}{images/Assistant_IA.png}
\caption{Diagramme de séquence — Fonctionnement de l’Assistant IA}
\label{fig:seq_assistant}
\end{figure}

Ce scénario relie la conversation, la validation humaine et l’exécution automatisée.

% ===============================================================
% 6. DIAGRAMME DE CLASSE
% ===============================================================

\section{Diagramme de classe}

Le diagramme de classe de figure~\ref{fig:class} décrit la structure statique d’Agent-Market. 


\figlead[0.78\textheight]{Le diagramme de classe ci-dessous synthétise ces relations assurent la cohérence des données multi-tenant et préparent l’implémentation des modules social, Ads et assistant.}
\begin{figure}[H]
\centering
\safeimage{width=1.05\textwidth,height=0.76\textheight}{images/class.png}
\caption{Diagramme de classe du système Agent-Market}
\label{fig:class}
\end{figure}

Le \textbf{Workspace} agrège utilisateurs, comptes, contenus, campagnes et artefacts d’agent. La classe \textbf{Compte} se spécialise en comptes publicitaires et comptes de réseaux sociaux ; les posts et campagnes publient respectivement sur ces comptes. Les relations de composition lient \textbf{AgentConversation} à \textbf{AgentMessage}, tandis qu’\textbf{AgentWorkflow} orchestre les traitements IA.

Le tableau~\ref{tab:classes} résume les classes du modèle ; la figure~\ref{fig:class} en donne la vue UML complète.

\begin{table}[H]
\centering
\caption{Dictionnaire des classes du système Agent-Market}
\label{tab:classes}
\renewcommand{\arraystretch}{1.55}
\setlength{\tabcolsep}{8pt}
\begin{tabular}{|p{0.26\textwidth}|p{0.33\textwidth}|p{0.33\textwidth}|}
\hline
\textbf{Classe} & \textbf{Description} & \textbf{Responsabilités principales} \\
\hline
Utilisateur & Compte applicatif & Authentification, profil, accès aux workspaces. \\
\hline
Workspace & Espace collaboratif & Membres, invitations, organisation des données. \\
\hline
Role & Permissions workspace & Contrôle d’accès. \\
\hline
Compte & Compte OAuth générique & Tokens, provider, rattachement workspace. \\
\hline
Comptes de réseaux sociaux & Spécialisation sociale & Publication et métriques sociales. \\
\hline
Comptes publicitaires & Spécialisation Ads & Campagnes et APIs publicitaires. \\
\hline
Post & Publication sociale & Brouillon, planification, multi-comptes. \\
\hline
Media & Fichier multimédia & Association aux posts/campagnes. \\
\hline
Campaign & Campagne marketing & Statut, plateforme, publication Ads. \\
\hline
LandingPage & Page de conversion & Collecte de leads. \\
\hline
Lead & Prospect & Attribution à une landing page. \\
\hline
AgentConversation & Dialogue IA & Historique et contexte. \\
\hline
AgentMessage & Message de conversation & Contenu utilisateur / assistant. \\
\hline
AgentWorkflow & Plan d’exécution IA & Création et suivi des étapes. \\
\hline
\end{tabular}
\renewcommand{\arraystretch}{1.0}
\end{table}

La figure~\ref{fig:class} met en évidence les multiplicités (workspace $1..*$ utilisateurs, posts liés à un média, campagnes liées aux comptes Ads, etc.) cohérentes avec le modèle MongoDB.


\section{Diagramme d’activité}

Les diagrammes d’activité formalisent les décisions et enchaînements du processus Assistant IA : d’abord conversation, planification et validation, puis exécution étape par étape jusqu’au résumé final.

\subsection{Conversation, planification et validation}

\figlead[0.88\textheight]{La figure~\ref{fig:act_plan} présente ce flux de conversation, planification et validation.}
\begin{figure}[H]
\centering
\safeimage{width=0.8\textwidth,height=0.6\textheight}{images/ai-assistant-activity-planification.png}
\caption{Diagramme d’activité — Conversation, planification et validation}
\label{fig:act_plan}
\end{figure}

Le processus part de l’ouverture de l’assistant et de la saisie d’un message. Après enregistrement de la conversation et planification asynchrone, le worker charge le contexte (workspace, comptes, outils) et analyse l’intent via le LLM. Deux chemins se distinguent : \textbf{conversation} (génération et affichage d’une réponse) ou \textbf{workflow} (construction et affichage du plan). La validation peut conduire au refus (annulation), à une approbation manuelle ou à un mode \textbf{auto}, avant le lancement asynchrone de l’exécution.

Cette activité complète la séquence Assistant IA en explicitant les décisions métier : nature de l’intent, nécessité d’un plan, mode de validation et bascule vers l’exécution. Elle montre aussi que le refus interrompt proprement le workflow, sans effet de bord sur les plateformes externes.



Ce processus place la validation humaine au centre du cycle, tout en permettant une exécution automatique lorsque le mode le prévoit.

\subsection{Exécution du workflow et résultats}

\figlead[0.88\textheight]{La figure~\ref{fig:act_exec} détaille cette boucle d’exécution et la gestion des erreurs.}
\begin{figure}[H]
\centering
\safeimage{width=0.8\textwidth,height=0.6\textheight}{images/ai-assistant-activity-execution.png}
\caption{Diagramme d’activité — Exécution du workflow et résultats}
\label{fig:act_exec}
\end{figure}

L’exécution décrit la boucle «~pour chaque étape~» : sélection de l’outil, puis action \textbf{externe} (connecteur Social/Ads) ou \textbf{interne} (création/mise à jour d’une ressource métier). En cas d’échec, une erreur critique marque le workflow en échec ; sinon l’échec est enregistré et le traitement poursuit. À la fin de la boucle, un résumé LLM est généré, le statut final est persisté et l’utilisateur est notifié.

Le point central est la distinction entre erreur critique (arrêt) et erreur non critique (poursuite), qui évite qu’un incident local n’annule inutilement l’ensemble du plan. La notification de progression et le résumé final rendent le traitement observable depuis l’interface.



Ces deux activités complètent les diagrammes de séquence de l’assistant en précisant les branchements métier (intent, validation, criticité des erreurs).

% ===============================================================
% 7. CONCLUSION
% ===============================================================

\section{Conclusion}

Ce chapitre a présenté la modélisation UML d’Agent-Market : cas d’utilisation (vue globale, publications, campagnes, assistant), diagrammes de séquence (authentification, OAuth, publication, assistant), diagramme de classe et diagrammes d’activité de planification puis d’exécution.

Ces modèles clarifient les interactions entre acteurs, modules et composants techniques. Le chapitre suivant détaille l’architecture déployée et les technologies retenues pour la réalisation.

\newpage

% ===============================================================
% CHAPTER 4
% ===============================================================

\refstepcounter{chapter}

\addcontentsline{toc}{chapter}{Chapitre \thechapter : Outils et technologies de développement}

\chapterpagetitle{Outils et technologies de développement}

% ===============================================================
% CHAPITRE — OUTILS ET TECHNOLOGIES
% ===============================================================

\section{Introduction}

Ce chapitre présente l’architecture déployée d’Agent-Market, puis les technologies et outils mobilisés pour sa réalisation. Nous détaillons d’abord la \textbf{vue système globale}, puis l’architecture de l’assistant IA en deux phases (\textbf{planification} et \textbf{exécution}), avant d’exposer la stack technique et l’environnement de travail.

\section{Architecture du système}

L’architecture d’Agent-Market se compose de plusieurs composants interconnectés, chacun ayant un rôle spécifique dans le traitement des demandes marketing, l’orchestration des tâches et la restitution des résultats. Voici une description détaillée de chaque composant et de leurs interactions à partir de la figure~\ref{fig:architecture_globale}.

\begin{figure}[H]
\centering
\safeimage{width=1.0\textwidth,height=0.5\textheight}{images/Architecture_Système_(Vue_Globale).png}
\caption{Architecture système d’Agent-Market — vue globale}
\label{fig:architecture_globale}
\end{figure}

\begin{itemize}[leftmargin=2em,itemsep=6pt]

\item \textbf{Utilisateur~:}
Le processus commence par l’utilisateur, marketeur qui interagit avec la plateforme via l’interface web pour gérer publications, campagnes, comptes connectés et assistant IA.

\textit{Interaction~:} l’utilisateur utilise l’interface (étape~1) pour consulter le dashboard, saisir des actions et suivre les notifications en temps réel.

\item \textbf{Interface web (Next.js / React)~:}
L’interface web est le point d’interaction entre l’utilisateur et l’application. Elle présente les écrans métier, envoie les requêtes authentifiées et affiche les mises à jour live.

\textit{Interaction~:} elle envoie des requêtes HTTP sécurisées par JWT à l’API métier (étape~2) et reçoit les événements via le canal temps réel (étape~10) pour actualiser l’UI sans rechargement complet.

\item \textbf{API métier (Fastify / Node.js)~:}
L’API métier gère la logique applicative, l’authentification JWT, la validation des entrées et l’orchestration entre persistance, file de tâches et événements.

\textit{Interaction~:} elle lit et enregistre les données dans MongoDB (étape~3), pousse les tâches via \texttt{LPUSH} vers Redis (étape~4), publie des événements (étape~8) et reçoit les retours HTTP internes du worker (étape~7).

\item \textbf{MongoDB~:}
MongoDB constitue la source de vérité documentaire (utilisateurs, workspaces, posts, campagnes, leads, conversations et workflows agent).

\textit{Interaction~:} l’API y persiste l’état métier ; le worker contribue indirectement via les mises à jour remontées à l’API.

\item \textbf{Redis (Celery + Pub/Sub)~:}
Redis sert de broker de messages~: file d’attente des tâches Celery et canal pub/sub pour le temps réel.

\textit{Interaction~:} il reçoit les tâches de l’API (étape~4), les livre au worker (étape~5), relaie les événements publiés (étapes~8--9) vers SocketCluster.

\item \textbf{Worker (Python / Celery)~:}
Le worker exécute les traitements longs~: génération IA, publication sociale, gestion Ads, imports et synchronisations.

\textit{Interaction~:} il consomme les tâches Redis (étape~5), appelle les services externes (étapes~6a--6c), puis notifie l’API en HTTP interne (étape~7).

\item \textbf{Temps réel (SocketCluster)~:}
SocketCluster assure la diffusion WebSocket des notifications vers le dashboard.

\textit{Interaction~:} il s’abonne au relais Redis pub/sub (étape~9) et pousse les mises à jour vers l’interface via le canal d’événements (étape~10).

\item \textbf{Services externes (IA, Social, Ads)~:}
Les fournisseurs IA (OpenAI / Anthropic, etc.), les réseaux sociaux (Meta, X, \ldots) et les plateformes Ads (Meta / Google) étendent les capacités métier hors du périmètre applicatif interne.

\textit{Interaction~:} le worker leur délègue la génération de contenu (6a), la publication (6b) et la gestion de campagnes (6c).

\end{itemize}

\subsection{Synthèse du flux global}

En résumé, l’utilisateur interagit via Next.js/React~; l’API Fastify orchestre JWT, MongoDB et Redis~; le worker Celery exécute les opérations lourdes vers l’IA, le social et les Ads~; Redis pub/sub et SocketCluster ramènent les événements vers l’interface. Cette chaîne garantit une séparation claire entre requête synchrone, traitement asynchrone et notification temps réel.

\section{Architecture de l’assistant IA — planification}

Au-delà de la vue système, l’assistant IA s’appuie sur la même chaîne technique pour une phase dédiée~: \textbf{comprendre / planifier}. La figure~\ref{fig:architecture_plan} détaille les composants et leurs interactions (A1~$\rightarrow$~A10).

\begin{figure}[H]
\centering
\safeimage{width=1.0\textwidth,height=0.5\textheight}{images/architecture_ai-assistant-planing.png}
\caption{Architecture de l’assistant IA — phase A (comprendre / planifier)}
\label{fig:architecture_plan}
\end{figure}

\begin{itemize}[leftmargin=2em,itemsep=6pt]

\item \textbf{Utilisateur~:}
Il exprime un besoin marketing en langage naturel dans le chat.

\textit{Interaction~:} envoi du message via le chat Next.js (A1).

\item \textbf{Chat (Next.js)~:}
Interface conversationnelle de l’assistant~; affiche ensuite le plan généré.

\textit{Interaction~:} appelle \texttt{POST /agent/chat} (A2), puis interroge le statut du job pour afficher le plan (A10).

\item \textbf{API (Fastify)~:}
Orchestre la création de la conversation/workflow et l’enqueue de la planification.

\textit{Interaction~:} persiste conversation et workflow dans MongoDB (A3), enfile \texttt{plan\_workflow} dans Redis (A4), reçoit le patch du planificateur (A9).

\item \textbf{MongoDB~:}
Stocke l’historique de conversation et l’état du workflow en cours de planification.

\textit{Interaction~:} lecture/écriture par l’API aux étapes A3 et lors des mises à jour ultérieures.

\item \textbf{Redis (bridge)~:}
File d’attente entre l’API et le worker pour la tâche de planification.

\textit{Interaction~:} reçoit \texttt{plan\_workflow} (A4) et la livre au worker (A5).

\item \textbf{Worker (Celery)~:}
Consomme la tâche asynchrone et démarre le module de planification.

\textit{Interaction~:} consomme la tâche (A5) puis lance le planificateur (A6).

\item \textbf{Planificateur~:}
Construit un plan multi-étapes à partir de l’intention utilisateur et du catalogue d’outils.

\textit{Interaction~:} injecte le catalogue (A7), interroge le LLM pour comprendre l’intention (A8), envoie un patch workflow à l’API (A9).

\item \textbf{Catalogue d’outils~:}
Référentiel des actions disponibles (social, Ads, ressources métier, etc.).

\textit{Interaction~:} fourni au planificateur pour borner le plan aux capacités réellement exécutables (A7).

\item \textbf{LLM~:}
Fournisseur d’intelligence artificielle pour la compréhension de la demande et la proposition de plan.

\textit{Interaction~:} reçoit le contexte du planificateur et retourne une intention structurée / un plan (A8).

\end{itemize}

\subsection{Synthèse de la phase A}

La phase A transforme un message utilisateur en plan validable~: chat $\rightarrow$ API $\rightarrow$ MongoDB/Redis $\rightarrow$ worker/planificateur/LLM, puis affichage du plan dans l’interface. Aucune action à effet de bord (publication, campagne) n’est encore exécutée.

\section{Architecture de l’assistant IA — exécution}

Une fois le plan validé, la phase \textbf{B — valider / exécuter} (B1~$\rightarrow$~B10) met en œuvre les étapes prévues. La figure~\ref{fig:architecture_exec} en précise les composants.

\begin{figure}[H]
\centering
\safeimage{width=1.0\textwidth,height=0.5\textheight}{images/architecture_ai-assistant-execution.png}
\caption{Architecture de l’assistant IA — phase B (valider / exécuter)}
\label{fig:architecture_exec}
\end{figure}

\begin{itemize}[leftmargin=2em,itemsep=6pt]

\item \textbf{Utilisateur~:}
Il examine le plan proposé et décide de l’approuver.

\textit{Interaction~:} validation du plan dans le chat (B1).

\item \textbf{Chat (Next.js)~:}
Déclenche l’exécution et affiche progression puis résumé.

\textit{Interaction~:} envoie \texttt{POST approve / execute} (B2) et reçoit le suivi via poll / WebSocket (B10).

\item \textbf{API (Fastify)~:}
Enfile l’exécution, persiste les mises à jour et expose le statut au chat.

\textit{Interaction~:} enqueue \texttt{execute\_workflow} dans Redis (B3), reçoit patches/événements (B8), met à jour MongoDB (B9), notifie l’UI (B10).

\item \textbf{Redis (bridge)~:}
Transporter la tâche d’exécution vers le worker.

\textit{Interaction~:} file \texttt{execute\_workflow} consommée en B4.

\item \textbf{Worker (Celery)~:}
Démarre le moteur d’exécution du workflow.

\textit{Interaction~:} consomme la tâche (B4) et démarre l’exécuteur (B5).

\item \textbf{Exécuteur~:}
Parcourt les étapes du plan dans l’ordre défini.

\textit{Interaction~:} dispatch chaque étape vers le dispatcher d’outils (B6) et remonte patches + events à l’API (B8).

\item \textbf{Dispatcher (outils)~:}
Route chaque étape vers le connecteur adapté.

\textit{Interaction~:} actions marketing vers Social \& Ads (B7a) et génération de contenu via LLM (B7b).

\item \textbf{Social \& Ads~:}
Plateformes externes de publication et de publicité.

\textit{Interaction~:} exécutent les actions demandées par le dispatcher (B7a) et retournent succès/erreur.

\item \textbf{LLM~:}
Génère les contenus nécessaires pendant l’exécution (textes, variantes, résumés intermédiaires selon les étapes).

\textit{Interaction~:} appelé par le dispatcher pour la génération (B7b).

\item \textbf{MongoDB~:}
Conserve l’état final et intermédiaire du workflow exécuté.

\textit{Interaction~:} mis à jour par l’API après chaque remontée d’événements (B9).

\end{itemize}

\subsection{Synthèse de la phase B}

Après validation humaine, l’API enfile l’exécution~; le worker active l’exécuteur et le dispatcher orchestre Social/Ads et LLM. Les événements remontent à l’API/MongoDB, puis le chat affiche le résumé. Planification et exécution restent ainsi découplées, tout en réutilisant la même infrastructure que l’architecture système globale.

% ===============================================================
\section{Technologies de la plateforme}

La plateforme repose sur les composants décrits ci-dessus. Les technologies suivantes correspondent à chaque couche~: interface web, API d’orchestration, worker asynchrone, moteur temps réel, stockage, fournisseurs IA et intégrations externes.

% Ligne techno : texte à gauche, logo à droite
% Contrôle global des tailles (modifiez ces valeurs) :
\newcommand{\techTextWidth}{0.78\textwidth}   % largeur du paragraphe
\newcommand{\techLogoColWidth}{0.18\textwidth}% largeur de la colonne logo
\newcommand{\techLogoHeight}{1.55cm}          % hauteur max du logo
\newcommand{\techLogoWidth}{0.95\linewidth}  % largeur max du logo dans sa colonne
% Usage :
%   \techrow{images/Gemini.png}{...}           % taille par défaut
%   \techrow[2.2cm]{images/Gemini.png}{...}    % hauteur personnalisée
\newcommand{\techrow}[3][\techLogoHeight]{%
  \par\noindent
  \begin{minipage}[c]{\techTextWidth}
  #3
  \end{minipage}\hfill
  \begin{minipage}[c]{\techLogoColWidth}
  \centering
  \safeimage{width=\techLogoWidth,height=#1,keepaspectratio}{#2}
  \end{minipage}\par
  \vspace{0.85em}%
}

\subsection{Frontend}

La couche frontend constitue l’interface utilisateur du dashboard marketing. Elle consomme l’API REST et reçoit les mises à jour en temps réel via WebSocket.

\subsubsection{Next.js}

\techrow[2.2cm]{images/Nextjs.png}{%
Next.js est un framework basé sur React permettant de développer des applications web performantes et optimisées. Il offre le rendu côté serveur (SSR), la génération statique (SSG) et un routage adapté au dashboard Agent-Market.%
}

\subsubsection{React}

\techrow[2.2cm]{images/React.png}{%
React est une bibliothèque JavaScript utilisée pour construire des interfaces interactives. Elle permet de créer des composants réutilisables et de gérer l’état de l’application de manière fluide.%
}

\subsubsection{CSS}

\techrow[2.2cm]{images/CSS.png}{%
CSS assure la mise en forme et la présentation visuelle de l’application~: styles, animations et mises en page responsives pour une expérience utilisateur cohérente.%
}

\subsection{Backend}

La couche backend regroupe l’API d’orchestration, le runtime Node.js, le worker asynchrone et le moteur de communication temps réel.

\subsubsection{Node.js}

\techrow[2.2cm]{images/node.png}{%
Node.js fournit l’environnement d’exécution JavaScript côté serveur. Il héberge l’API Fastify et permet de traiter efficacement les requêtes concurrentes du dashboard.%
}

\subsubsection{Fastify}

\techrow[2.2cm]{images/Fastify.png}{%
Fastify est un framework backend Node.js reconnu pour sa rapidité et sa faible empreinte. Il structure l’API métier (JWT, validation, plugins) de façon modulaire et performante.%
}

\subsubsection{Python}

\techrow[2.2cm]{images/Python.png}{%
Python est utilisé pour le worker d’automatisation, les connecteurs externes et les appels aux fournisseurs IA. Il convient aux traitements longs et aux intégrations data/IA.%
}

\subsubsection{Celery}

\techrow[2.2cm]{images/Celery.png}{%
Celery orchestre les tâches asynchrones Python~: publication sociale, synchronisations Ads, planification de posts et exécution des workflows agent sans bloquer l’API.%
}

\subsubsection{SocketCluster}

\techrow[2.2cm]{images/SocketCluster.png}{%
SocketCluster assure la communication temps réel (WebSocket). Il relaie les événements Redis vers le dashboard pour les notifications live (publications, workflows, sync).%
}

\subsection{Base de données et messagerie}

Les couches de persistance et de messagerie assurent le stockage des données métier, la gestion des files de tâches et la diffusion des événements applicatifs.

\subsubsection{MongoDB}

\techrow[2.2cm]{images/MongoDB.png}{%
MongoDB est la base documentaire NoSQL de la plateforme. Elle stocke utilisateurs, workspaces, posts, campagnes, leads, conversations et workflows au format flexible JSON/BSON.%
}

\subsubsection{Redis}

\techrow[2.2cm]{images/Redis.png}{%
Redis sert de broker Celery et de canal pub/sub. Il enfile les tâches API$\rightarrow$worker et relaie les événements temps réel vers SocketCluster.%
}

\subsection{Fournisseurs d’intelligence artificielle}

L’assistant et la génération de contenu s’appuient sur plusieurs fournisseurs LLM configurables par workspace.

\subsubsection{Google Gemini}

\techrow[2.2cm]{images/Gemini.png}{%
Gemini est un fournisseur LLM de Google utilisé pour la compréhension des demandes, la planification de workflows et la génération de contenus marketing (texte et, selon les cas, multimédia).%
}

\subsubsection{OpenAI}

\techrow[2.2cm]{images/chatgpt.jpg}{%
OpenAI fournit des modèles de langage pour la génération de contenu, le résumé de résultats et l’assistance conversationnelle au sein de l’agent.%
}

\subsection{Intégrations sociales et publicitaires}

Les connecteurs OAuth et API relient Agent-Market aux plateformes de diffusion et d’acquisition.

\subsubsection{Meta}

\techrow[2.2cm]{images/Meta.png}{%
Meta (Facebook / Instagram et Meta Ads) permet la connexion OAuth des pages, la publication sociale et la gestion/synchronisation des campagnes publicitaires.%
}

\subsubsection{Google Ads}

\techrow[2.2cm]{images/google_ads.png}{%
Google Ads est intégré pour connecter les comptes publicitaires, publier ou synchroniser les campagnes et récupérer les performances.%
}


\section{Environnement et outils de développement}

Au-delà de la stack applicative, l’environnement de travail garantit le déploiement cohérent et reproductible de l’ensemble des services.

\subsection{Docker Compose}

\techrow[2.2cm]{images/DockerCompose.png}{%
Docker Compose orchestre les conteneurs (frontend, API, MongoDB, Redis, worker Celery, SocketCluster, etc.) dans une configuration unique, reproductible en développement comme en production.%
}

\section{Conclusion}

Ce chapitre a présenté l’architecture système et agent, puis la stack technologique d’Agent-Market~: frontend, backend, persistance, fournisseurs IA (Gemini, OpenAI, Anthropic), intégrations Meta / Google Ads / X, et l’environnement Docker Compose.

Le choix de ces technologies a été guidé par des critères de performance, de scalabilité et de maintenabilité. Elles constituent la base technique du projet et permettent de garantir une architecture moderne, robuste et évolutive.

Dans le chapitre suivant, nous présenterons la réalisation de la solution à travers les interfaces de l’application et leurs explications.

\newpage

% ===============================================================
% CHAPTER 5 — RÉALISATION
% ===============================================================

\refstepcounter{chapter}

\addcontentsline{toc}{chapter}{Chapitre \thechapter : Réalisation}

\chapterpagetitle{Réalisation}

% ===============================================================
% CHAPITRE — RÉALISATION
% ===============================================================

\section{Introduction}

Dans cette phase de réalisation, nous présentons l’implémentation de la plateforme \textbf{Agent Market} à travers ses principales fonctionnalités, ses interfaces graphiques et les parcours utilisateur couvrant l’authentification, la gestion des réseaux sociaux, les campagnes publicitaires, la génération de leads et l’assistant IA.

L’application authentifiée s’articule autour d’une barre latérale organisée en trois sections — \textbf{Workspace}, \textbf{Marketing} et \textbf{Operations} —, d’une barre supérieure affichant le workspace actif, ainsi que d’un bouton flottant permettant d’accéder à l’assistant IA depuis n’importe quelle page.


\section{Interfaces d’authentification et d’accueil}

\subsection{Page de connexion}

Cette interface constitue le point d’entrée de la plateforme. Elle présente le logo Agent Market, le titre de l’application et le sous-titre \textit{Sign in to your account}. L’utilisateur saisit son adresse e-mail et son mot de passe dans un formulaire centré, puis valide via le bouton \textbf{Sign in}.

\begin{figure}[H]
\centering
\safeimage{width=0.6\textwidth,height=0.4\textheight}{images/login.png}
\caption{Interface de connexion Agent Market}
\label{fig:ui_login}
\end{figure}

En cas d’échec, un message d’erreur indique que les identifiants sont invalides. Un lien permet également de basculer vers la page d’inscription pour les nouveaux utilisateurs.

\subsection{Page d’inscription}

Cette interface permet de créer un compte sur la plateforme. Elle propose les champs \textbf{Full Name}, \textbf{Email}, \textbf{Password} et \textbf{Confirm Password}, ainsi que le bouton \textbf{Create Account}. Des contrôles côté interface vérifient que tous les champs sont renseignés, que les mots de passe correspondent et que le mot de passe respecte une longueur minimale. Dans le cas d’une invitation à rejoindre un workspace, un message spécifique indique que le compte permettra de rejoindre l’espace de travail invité.

\begin{figure}[H]
\centering
\safeimage{width=0.6\textwidth,height=0.6\textheight}{images/register.png}
\caption{Interface d’inscription}
\label{fig:ui_register}
\end{figure}

\subsection{Tableau de bord Overview}

Après authentification, l’utilisateur accède au tableau de bord \textbf{Overview}, point d’entrée principal de l’espace de travail authentifié. Cette page offre une vue synthétique et immédiatement lisible de l’activité du workspace, destinée à permettre un diagnostic rapide de l’état éditorial sans naviguer vers des écrans plus détaillés. En tête de page, quatre indicateurs clés résument le volume et la santé des publications : \textbf{Total Posts} (ensemble des contenus créés), \textbf{Scheduled} (publications planifiées en attente), \textbf{Published} (contenus déjà diffusés avec succès) et \textbf{Failed} (échecs de publication à traiter). Sous ces métriques, la liste des publications récentes présente un aperçu chronologique des derniers contenus, avec leur statut et les informations essentielles pour un suivi au quotidien. Une zone dédiée aux comptes connectés rappelle les pages et profils déjà rattachés au workspace, facilitant le contrôle de la couverture multi-plateforme. Un bouton \textbf{+ New Post} permet de démarrer rapidement une nouvelle publication depuis le tableau de bord. Lorsque aucune donnée n’est encore disponible — workspace nouvellement créé ou absence de contenus —, des états vides guident explicitement l’utilisateur vers les actions prioritaires : création d’un premier post ou connexion d’un compte via \textbf{Connect Account}, afin de rendre la prise en main progressive et sans ambiguïté.

\begin{figure}[H]
\centering
\safeimage{width=1.0\textwidth,height=0.45\textheight}{images/overview.png}
\caption{Tableau de bord Overview}
\label{fig:ui_overview}
\end{figure}

\section{Interfaces de gestion sociale}

\subsection{Connexion des comptes sociaux et publicitaires}

L’écran \textbf{Accounts} centralise la connexion et la gestion des pages et profils sociaux destinés à la publication, ainsi que des comptes publicitaires Google Ads et Meta Ads. Il constitue le prérequis opérationnel pour la diffusion de contenus et le pilotage des campagnes : sans compte rattaché, les parcours de publication et de publicité ne peuvent aboutir. Deux sections distinctes structurent l’interface afin de séparer clairement les usages éditoriaux et publicitaires : \textbf{Social Accounts}, accessible via le bouton \textbf{Add Account}, regroupe les profils et pages des réseaux sociaux ; \textbf{Ads Accounts}, accessible via \textbf{Connect Ads Account}, regroupe les comptes publicitaires.

\begin{figure}[H]
\centering
\safeimage{width=1.0\textwidth,height=0.45\textheight}{images/accounts.png}
\caption{Interface de gestion des comptes connectés}
\label{fig:ui_accounts}
\end{figure}

Chaque compte affiche son statut (\textbf{Active} lorsque les jetons sont valides, \textbf{Unauthorized} en cas d’expiration ou de révocation des autorisations) ainsi que des actions de synchronisation ou de rafraîchissement permettant de maintenir à jour les métadonnées et les permissions. Des fenêtres modales guident le choix de la plateforme cible — Facebook, Instagram, TikTok, LinkedIn, X pour le social, Google Ads et Meta Ads pour la publicité — avant de lancer le flux OAuth vers le fournisseur externe. Après authentification réussie chez le tiers, une page de sélection d’entités permet à l’utilisateur de choisir précisément la page, le profil ou le compte client à rattacher au workspace, évitant ainsi d’importer automatiquement l’ensemble des ressources disponibles et garantissant un contrôle fin du périmètre connecté.

\subsection{Liste des publications}

L’interface \textbf{Posts} constitue le centre de gestion de l’ensemble des publications du workspace. Elle permet de suivre le cycle de vie des contenus — de la rédaction à la diffusion — et d’intervenir rapidement sur les éléments à corriger ou à republier. Un bouton \textbf{New Post} ouvre le parcours de création, tandis que des onglets de filtre (\textbf{All}, \textbf{Draft}, \textbf{Scheduled}, \textbf{Published}, \textbf{Failed}) facilitent le basculement entre les différentes étapes du pipeline éditorial. Le tableau principal présente, pour chaque publication, le contenu (texte et aperçu média), le statut courant, la date planifiée le cas échéant, ainsi que les comptes sociaux associés à la diffusion.

\begin{figure}[H]
\centering
\safeimage{width=1.0\textwidth,height=0.5\textheight}{images/posts.png}
\caption{Interface de liste des publications}
\label{fig:ui_posts}
\end{figure}

\subsection{Création et édition d’une publication}

La création d’une publication s’effectue via un assistant multi-étapes structuré en quatre phases (\textbf{Type}, \textbf{Content}, \textbf{Media}, \textbf{Preview}), conçu pour guider l’utilisateur de manière progressive et réduire les erreurs de configuration. À l’étape \textbf{Type}, l’utilisateur choisit le format du contenu — \textbf{Image post} ou \textbf{Video post} — selon le type de média à diffuser. L’étape \textbf{Content} permet de rédiger manuellement le texte ou de le générer à l’aide de l’IA en précisant un objectif, un ton rédactionnel et d’éventuels hashtags, afin d’adapter le message au contexte marketing. L’étape \textbf{Media} offre ensuite la possibilité de joindre des fichiers existants ou de générer des visuels via l’IA, avant que l’étape \textbf{Preview} ne présente un aperçu du rendu par plateforme, utile pour vérifier la cohérence du contenu sur chaque réseau. L’écran d’édition complète ce parcours en permettant de sélectionner les comptes destinataires, de planifier une date et une heure de publication, d’ajouter un premier commentaire (Facebook / Instagram) et de définir un ciblage géographique lorsque la plateforme le permet. 
\begin{figure}[H]
\centering
\safeimage{width=1.0\textwidth,height=0.4\textheight}{images/post_wizard.png}
\caption{Assistant de création de publication}
\label{fig:ui_post_wizard}
\end{figure}

\begin{figure}[H]
\centering
\safeimage{width=1.0\textwidth,height=0.4\textheight}{images/post_editor.png}
\caption{Éditeur de publication avec prévisualisation}
\label{fig:ui_post_editor}
\end{figure}

Les actions finales proposées — \textbf{Save Draft}, \textbf{Schedule} et \textbf{Publish} — couvrent respectivement la sauvegarde pour reprise ultérieure, la planification différée et la publication immédiate.

\subsection{Calendrier éditorial}

L’interface \textbf{Calendar} affiche une grille mensuelle des publications, offrant une vision temporelle de la planification éditoriale du workspace. Contrairement à la liste \textbf{Posts}, qui privilégie le détail opérationnel de chaque contenu, le calendrier met l’accent sur la répartition dans le temps et sur la densité éditoriale jour après jour.

Les publications planifiées peuvent en outre être déplacées par glisser-déposer afin d’être reprogrammées sans quitter le calendrier~: cette interaction directe facilite l’ajustement du planning lorsque les priorités marketing évoluent, qu’un créneau se libère ou qu’un conflit de publication apparaît.

\begin{figure}[H]
\centering
\safeimage{width=1.0\textwidth,height=0.5\textheight}{images/calendar.png}
\caption{Calendrier éditorial}
\label{fig:ui_calendar}
\end{figure}

Une légende visuelle distingue clairement les statuts \textbf{Draft}, \textbf{Scheduled}, \textbf{Published} et \textbf{Failed}, permettant d’identifier d’un coup d’oeil les contenus encore en préparation, ceux prêts à paraître, ceux déjà diffusés et ceux ayant échoué. Cette lecture colorée et synthétique fluidifie le pilotage de la planification éditoriale et aide l’équipe à détecter rapidement les déséquilibres ou les actions à entreprendre.


\subsection{Médiathèque}

La page \textbf{Media Library} centralise la gestion des images et vidéos du workspace, servant de réservoir de ressources réutilisables lors de la création de publications ou de contenus marketing. L’utilisateur peut importer de nouveaux fichiers via le bouton \textbf{Upload}, sélectionner plusieurs éléments pour une suppression groupée afin d’alléger la bibliothèque, et parcourir la grille média paginée pour retrouver rapidement un actif. Chaque élément est présenté sous forme de vignette, ce qui facilite l’identification visuelle des médias disponibles. Un état vide invite explicitement à téléverser les premiers fichiers lorsque la bibliothèque est encore vide, guidant ainsi l’utilisateur vers la première étape d’alimentation de son espace média.

\begin{figure}[H]
\centering
\safeimage{width=1.0\textwidth,height=0.5\textheight}{images/media.png}
\caption{Interface de la médiathèque}
\label{fig:ui_media}
\end{figure}

\subsection{Rapports de performance sociale}

L’écran \textbf{Reports} consolide les analytics par plateforme sociale afin d’offrir une lecture unifiée de la performance des comptes connectés. L’utilisateur sélectionne d’abord un compte, définit une période d’analyse, puis synchronise les métriques depuis la plateforme source pour disposer de données à jour. Il consulte ensuite des cartes de statistiques synthétiques et des graphiques d’évolution — notamment sur l’engagement et les followers — qui permettent de comparer les périodes et d’identifier les tendances. Des messages d’état vides indiquent clairement lorsqu’aucun compte n’est connecté ou lorsqu’aucune métrique n’a encore été importée, orientant l’utilisateur vers la connexion d’un compte ou le lancement d’une première synchronisation avant toute analyse.

\begin{figure}[H]
\centering
\safeimage{width=1.0\textwidth,height=0.5\textheight}{images/reports.png}
\caption{Interface des rapports sociaux}
\label{fig:ui_reports}
\end{figure}

\section{Interfaces marketing et publicitaires}

\subsection{Hub des campagnes Ads}

La page \textbf{Ads} constitue le point d’entrée de la gestion publicitaire Google Ads et Meta Ads au sein du workspace. Elle offre une vue d’ensemble destinée à évaluer rapidement la santé des investissements publicitaires sans naviguer immédiatement vers les écrans détaillés. En tête de page, des indicateurs globaux résument le budget alloué, les dépenses réalisées, le taux de clics (CTR) et le nombre de campagnes actives. Des panneaux par plateforme, munis d’actions de synchronisation, permettent de rafraîchir les données issues de Google Ads et Meta Ads. Un tableau \textbf{All Campaigns} regroupe ensuite l’ensemble des campagnes pour un suivi transversal. Des raccourcis facilitent enfin l’accès à la performance détaillée, à la gestion des comptes publicitaires et à la création d’une nouvelle campagne via \textbf{New Campaign}.

\begin{figure}[H]
\centering
\safeimage{width=1.0\textwidth,height=0.5\textheight}{images/ads.png}
\caption{Hub de gestion des campagnes publicitaires}
\label{fig:ui_ads}
\end{figure}

\subsection{Performance publicitaire}

L’écran \textbf{Ads Performance} permet d’analyser les insights d’un compte Google Ads ou Meta Ads. Après sélection du compte, l’utilisateur charge les rapports (\textbf{Run report} / \textbf{Load insights}) afin de consulter les indicateurs de performance. Lorsque aucun compte publicitaire n’est connecté, un message guide l’utilisateur vers la page Accounts.

\begin{figure}[H]
\centering
\safeimage{width=0.95\textwidth}{images/ads_performance.png}
\caption{Interface de performance publicitaire}
\label{fig:ui_ads_performance}
\end{figure}

\subsection{Landing pages}

L’interface \textbf{Landing Pages} permet de créer et publier des pages de conversion. Chaque carte affiche le titre, le slug public (\texttt{/lp/...}), le statut, les visites et les leads collectés. Les actions proposées incluent \textbf{Copy Link}, \textbf{Publish} / \textbf{Unpublish}, \textbf{Edit} et la suppression. Un assistant de création guide l’utilisateur de la description du besoin jusqu’à la génération assistée par IA, l’édition du contenu (titre, accroche, corps, CTA, champs du formulaire) et la publication.

\begin{figure}[H]
\centering
\safeimage{width=0.95\textwidth}{images/landing_pages.png}
\caption{Interface de gestion des landing pages}
\label{fig:ui_landing_pages}
\end{figure}

\subsection{Collecte des leads}

La page \textbf{Leads} centralise les prospects capturés via les formulaires des landing pages. Elle affiche un compteur de leads, un bouton \textbf{Export CSV}, et un tableau regroupant la date, les champs dynamiques, la page source et les paramètres UTM. Un état vide indique qu’aucune conversion n’a encore été enregistrée et invite à publier une landing page pour démarrer la collecte.

\section{Interfaces de l’Assistant IA}

\subsection{Panneau conversationnel AI Agent}

L’assistant IA n’est pas une page isolée : il est accessible depuis toutes les écrans authentifiés via un bouton flottant. Le panneau latéral affiche l’en-tête \textbf{AI Agent}, un historique des conversations, un bouton de nouvelle session et une zone de chat. L’utilisateur formule sa demande en langage naturel (\textit{Describe what you want to accomplish…}) et peut joindre des pièces. Lorsque le chat est vide, des suggestions concrètes sont proposées (création de post Facebook, campagne Google Search, brouillon Meta Ads).

\begin{figure}[H]
\centering
\safeimage{width=0.75\textwidth}{images/agent_panel.png}
\caption{Panneau de l’Assistant IA}
\label{fig:ui_agent_panel}
\end{figure}

\subsection{Planification, validation et exécution d’un workflow}

Après analyse de la demande, l’assistant génère un workflow structuré présenté sous forme de carte. L’utilisateur peut \textbf{Approve \& run}, rejeter le plan, ou arrêter une exécution en cours. Une bannière \textit{Working…} indique le traitement en cours, tandis que le suivi des étapes est restitué dans le fil de conversation. Ce mécanisme garantit que les actions marketing automatisées (publication, création de campagne, génération de contenu) ne sont lancées qu’après validation explicite de l’utilisateur.

\begin{figure}[H]
\centering
\safeimage{width=0.75\textwidth}{images/agent_workflow.png}
\caption{Carte de workflow IA à valider}
\label{fig:ui_agent_workflow}
\end{figure}

\subsection{Configuration des intégrations IA}

Réservée aux administrateurs du workspace, la page \textbf{AI Integrations} permet de configurer les fournisseurs (Gemini, OpenAI, Anthropic), de tester la connexion, d’activer les fonctionnalités IA et de paramétrer l’assistant marketing. Un historique des exécutions IA complète cette interface, offrant une traçabilité des appels effectués depuis le workspace.

\begin{figure}[H]
\centering
\safeimage{width=0.95\textwidth}{images/ai_integrations.png}
\caption{Interface de configuration des intégrations IA}
\label{fig:ui_ai_integrations}
\end{figure}

\section{Interfaces d’administration du workspace}

\subsection{Gestion d’équipe}

La page \textbf{Team} permet aux administrateurs de renommer le workspace, d’inviter des membres (e-mail, rôle Member / Admin) et de gérer les invitations en attente. La liste des membres affiche le rôle de chacun et propose, pour les administrateurs, une action de retrait. Les membres non administrateurs consultent uniquement la composition de l’équipe et leurs badges de rôle.

\begin{figure}[H]
\centering
\safeimage{width=0.95\textwidth}{images/team.png}
\caption{Interface de gestion d’équipe}
\label{fig:ui_team}
\end{figure}

\subsection{Profil et préférences}

L’écran \textbf{Profile} permet de modifier les informations personnelles (nom, e-mail) et le mot de passe. La page \textbf{Preferences}, réservée aux administrateurs, centralise les réglages régionaux (fuseau horaire, format de date/heure), le comportement de l’agent IA (auto-approbation des workflows), les defaults de publication et les options liées à l’analyse des commentaires. Une prévisualisation en direct accompagne la sauvegarde des préférences.

\begin{figure}[H]
\centering
\safeimage{width=0.95\textwidth}{images/preferences.png}
\caption{Interface des préférences du workspace}
\label{fig:ui_preferences}
\end{figure}

\section{Conclusion}

Dans cette phase de réalisation, nous avons concrétisé la plateforme Agent Market en mettant en œuvre les interfaces conçues lors de l’analyse et de la conception. Les modules d’authentification, de gestion sociale multi-canaux, de campagnes publicitaires, de landing pages, de leads et d’assistant IA sont désormais opérationnels au sein d’une interface unifiée. Cette étape a été le fruit d’un travail méthodique visant à offrir une expérience utilisateur cohérente, guidée et adaptée aux besoins marketing multi-plateformes de Tripass.

\newpage

% ===============================================================
% CONCLUSION GÉNÉRALE ET PERSPECTIVES
% ===============================================================
\centeredgreentitle{Conclusion Générale et Perspectives}
\addfrontmattertotoc{Conclusion Générale et Perspectives}

En conclusion de ce stage chez \textbf{Tripass}, l’expérience acquise a été enrichissante et formatrice, offrant bien plus qu’une simple immersion dans le développement web full-stack. J’ai eu l’opportunité de concevoir et de développer la plateforme \textbf{Agent Market}, une solution intelligente de gestion marketing multi-plateformes visant à centraliser les publications sociales, les campagnes publicitaires, la génération de leads et l’automatisation assistée par IA.

La réalisation d’Agent Market a permis de mettre en place une architecture distribuée composée d’une interface Next.js/React, d’une API Fastify, d’un worker Python/Celery, d’un moteur temps réel SocketCluster, ainsi que de MongoDB et Redis pour la persistance et la messagerie. Ce projet m’a permis de renforcer mes compétences en développement frontend et backend, en intégration d’APIs RESTful, en authentification JWT/OAuth, et en orchestration de tâches asynchrones. La connexion aux réseaux sociaux (LinkedIn, Meta, TikTok, X) et aux plateformes publicitaires (Google Ads, Meta Ads) a illustré l’importance de la synchronisation des données et de l’automatisation des processus dans un environnement marketing multi-canal.

Par ailleurs, le développement de l’assistant IA a introduit une dimension conversationnelle au cœur de la plateforme : l’utilisateur exprime un besoin en langage naturel, le système planifie un workflow, le valide auprès de l’utilisateur, puis l’exécute via le worker. L’intégration de fournisseurs LLM (Gemini, OpenAI, Anthropic), la génération de contenu, la création assistée de campagnes et de landing pages m’ont permis d’approfondir mes compétences en intelligence artificielle appliquée et de comprendre les défis de l’intégration d’agents intelligents dans une application métier réelle.

Au-delà des aspects techniques, ce stage m’a permis de saisir les complexités de la gestion de projet (méthodologie Waterfall), de comprendre les dynamiques d’une équipe pluridisciplinaire et d’apprécier l’importance de l’innovation continue. Les interactions quotidiennes avec des professionnels expérimentés — notamment mon encadrant pédagogique Mme. Mariem Malaine, Mr. Abid Khirani et Mr. Amine Ahbouch — m’ont permis de développer des compétences transversales, allant de la résolution de problèmes techniques à la communication efficace.

Les défis rencontrés — intégrations OAuth multi-plateformes, fiabilité des tâches asynchrones, cohérence des données entre systèmes externes et base locale, expérience utilisateur unifiée — ont été autant de sources d’apprentissage, renforçant ma capacité à m’adapter rapidement et à trouver des solutions pragmatiques. Cette expérience m’a également confronté à l’importance cruciale de la collaboration et du partage de connaissances au sein d’un environnement de travail dynamique.

En somme, ce stage chez Tripass a été une étape décisive dans mon parcours professionnel. Il a consolidé mes compétences techniques, élargi ma vision du développement d’applications distribuées et de l’IA appliquée au marketing digital, et m’a offert des perspectives enrichissantes quant aux opportunités et aux défis du monde professionnel.

\newpage

% ===============================================================
% RÉFÉRENCES
% ===============================================================
\centeredgreentitle{Références}
\addfrontmattertotoc{Références}

\begin{enumerate}[leftmargin=2em,itemsep=6pt,label={[\arabic*]}]
\item Vercel Inc., << Next.js Documentation >>, Documentation officielle Next.js, 2024, \url{https://nextjs.org/docs}, consulté le 26/07/2026.

\item Meta Platforms Inc., << React Documentation >>, Documentation officielle React, 2024, \url{https://react.dev/learn}, consulté le 26/07/2026.

\item Matteo Collina et al., << Fastify Documentation >>, Documentation officielle Fastify, 2024, \url{https://fastify.dev/docs/latest/}, consulté le 26/07/2026.

\item MongoDB Inc., << MongoDB Manual >>, Documentation officielle MongoDB, 2024, \url{https://www.mongodb.com/docs/}, consulté le 26/07/2026.

\item Redis Ltd., << Redis Documentation >>, Documentation officielle Redis, 2024, \url{https://redis.io/docs/}, consulté le 26/07/2026.

\item Ask Solem et al., << Celery -- Distributed Task Queue >>, Documentation officielle Celery, 2024, \url{https://docs.celeryq.dev/}, consulté le 26/07/2026.

\item OpenAI, << OpenAI API Documentation >>, Documentation officielle OpenAI, 2024, \url{https://platform.openai.com/docs}, consulté le 26/07/2026.

\item Google, << Gemini API Documentation >>, Google AI for Developers, 2024, \url{https://ai.google.dev/gemini-api/docs}, consulté le 26/07/2026.

\item Anthropic, << Claude API Documentation >>, Documentation officielle Anthropic, 2024, \url{https://docs.anthropic.com/}, consulté le 26/07/2026.

\item Meta Platforms Inc., << Meta Marketing API >>, Documentation développeurs Meta, 2024, \url{https://developers.facebook.com/docs/marketing-apis/}, consulté le 26/07/2026.

\item Google, << Google Ads API Documentation >>, Google Developers, 2024, \url{https://developers.google.com/google-ads/api/docs/start}, consulté le 26/07/2026.

\item Object Management Group, << Unified Modeling Language (UML) >>, Spécification OMG, 2017, \url{https://www.omg.org/spec/UML/}, consulté le 26/07/2026.

\item Winston W. Royce, << Managing the Development of Large Software Systems >>, Proceedings of IEEE WESCON, 1970, pp.~1--9.

\item Docker Inc., << Docker Compose Documentation >>, Documentation officielle Docker, 2024, \url{https://docs.docker.com/compose/}, consulté le 26/07/2026.

\item SocketCluster, << SocketCluster Documentation >>, Documentation officielle SocketCluster, 2024, \url{https://socketcluster.io/}, consulté le 26/07/2026.
\end{enumerate}

\end{document}
